% Condensed appendix for main.tex (v3). The complete appendix is appendix_extended.tex, compiled by extended_appendix.tex
% into extended_appendix.pdf for the supplement. Every number here is copied from it; tables in tables/ are generated by
% scripts/make_tables.py from papers/analysis/*.json.

\section*{Appendix overview}
Each section below supports the main-text sections listed. The extended appendix in the supplementary material contains every analysis and the full ledger, including exploratory work that is not part of the paper's claims (preference-training runs, a constraint-problem pilot, code experiments, and composition studies).

\begin{center}
\footnotesize
\setlength{\tabcolsep}{4pt}
\renewcommand{\arraystretch}{1.08}
\begin{tabular}{cL{10.6cm}c}
\toprule
& Contents & Main text \\
\midrule
\ref{app:setup} & Pools, prompts and the conditions each is used in, why the default rubrics invite comparison, the Figure~1 example & \ref{sec:intro}, \ref{sec:setup} \\
\ref{app:survey} & Graders, libraries, RL pipelines, and papers that grade against one reference, and what their prompts say & \ref{sec:intro} \\
\ref{app:prevalence} & How often questions admit several answers in nine datasets (LLM labels) & \ref{sec:cost}, \ref{sec:fix} \\
\ref{app:ledger} & Verdict totals and every registered prediction that was not confirmed & all \\
\ref{app:mace} & \textsc{MACE} domains, answer-list completeness, and answer matching & \ref{sec:setup}, \ref{sec:fix} \\
\ref{app:grid} & Answer labelings, the label audit, and grading under each labeling & \ref{sec:setup}, \ref{sec:p2} \\
\ref{app:estimator} & Exact scorer, bare-answer reference on MATH-500, and selection & \ref{sec:setup}, \ref{sec:p2} \\
\ref{app:reasoning} & Reasoning and thinking judges & \ref{sec:setup}, \ref{sec:p1} \\
\ref{app:rotation} & Swap test: all cells, the MT-Bench template, and the pool-level plot & \ref{sec:p1} \\
\ref{app:mechanism} & How the form of the reference (expression, masked solution) changes its pull & \ref{sec:p1} \\
\ref{app:frontier} & Frontier API judges on 16-candidate subsamples & \ref{sec:p2} \\
\ref{app:acceptance} & Blind acceptance and a second blind call, the full-credit sentence, benchmark scores with our rubric, and a second reference draw & \ref{sec:p1}, \ref{sec:p3}, \ref{sec:cost} \\
\ref{app:r24} & Benchmark graders run verbatim, and the full answer list & \ref{sec:cost}, \ref{sec:fix} \\
\ref{app:r25} & Judges that reason before deciding & \ref{sec:cost} \\
\ref{app:r26} & Direct prompt, edited HLE judge, grading with many wrong answers, and open RL-reward verifiers & \ref{sec:cost}, \ref{sec:fix} \\
\ref{app:r27} & P3 under correctness prompts, and hard single-answer questions & \ref{sec:p3}, \ref{sec:fix} \\
\ref{app:r30} & The swap under a prompt that does not instruct comparison, and a chosen reference & \ref{sec:p1}, \ref{sec:cost}, \ref{sec:fix} \\
\ref{app:r31} & The swap under a prompt that forbids using the reference & \ref{sec:p1}, \ref{sec:fix} \\
\ref{app:mitigation} & Blind-then-reveal, the max rule (post hoc), and other mitigations & \ref{sec:p2}, \ref{sec:fix} \\
\bottomrule
\end{tabular}
\end{center}

%==========================================================================================
\FloatBarrier
\section{Setup details}
\label{app:setup}

Every number in this condensed appendix also appears in the extended appendix. The ledger keeps the names used in the registrations, calling the swap test \emph{rot\-ation} and the margin the \emph{an\-choring margin}; the two rubric wordings are $a$ and $b$.

\paragraph{Pools.}
The pools are summarized in \cref{sec:setup}, and the extended appendix tabulates each with its generator and judges. The benchmark answers are one answer per model from Llama-3.1-8B-Instruct, Haiku~4.5, Sonnet~5, and Opus~5 to 149 \textsc{MACE} questions with several correct answers (596 answers). \paragraph{Scoring.}
Open judges run in vLLM at temperature 0 and produce at most four tokens. We read the top-20 log-probabilities of the first generated token and take the probability-weighted mean over the score tokens 0 to 10, renormalized over their captured probability. Qwen2.5, gemma-2, and Mistral write ``10'' as the two tokens ``1'' and ``0''; the original scorer assigned all probability on ``1'' to 10 or to 1 according to the greedy reply, and the pre-specified exact scorer of \cref{app:estimator} runs one more forward step to split it. Mistral-7B's tokenizer has no single-token digits, so its scores are the integers parsed from its greedy replies. Sonnet~5 through the Anthropic API ran with extended thinking disabled and returns its score through a forced, schema-validated tool call; GPT-4.1 uses the same expected-score estimator as the open judges, but its API returns at most five alternatives per token.

\paragraph{Measures.}
For each task we select the sample with the highest rubric-$a$ score; when several samples tie at the maximum, accuracy is the mean validity of the tied samples. False rejection uses rubric-$a$ scores of oracle-valid samples whose extracted answer (the last \texttt{\textbackslash boxed\{\}} content, loosely normalized, or the whole reply if there is none) differs from the reference. Both are averaged per task and bootstrapped over tasks (1,000 resamples). The anchoring \emph{margin} is the mean score of valid samples that match the reference minus the mean score of valid samples that do not. \emph{Inflation}, the endpoint of the early registered experiments, is the standardized rubric-$b$ score of the selected sample minus the mean standardized rubric-$b$ score of samples with the same validity, against a permutation null, on tasks with pass rate in $[0.2, 0.8]$, and reported as the referenced-minus-blind difference per task; the extended appendix gives the full definition.

\paragraph{Score rubrics and the conditions that use them.}
The judge receives a system message with the rubric and a user message containing ``Problem:'', optionally ``Reference answer (trusted):'', and ``Candidate solution:'' followed by ``Score:''. \emph{Blind condition:}
\begin{quote}\small
(a) You are a strict grader. Read the question and the candidate response. Rate how correct and appropriate the response is on a scale from 0 (completely wrong or off-task) to 10 (fully correct). Judge the content, not the writing. Respond with only the integer score.\\[2pt]
(b) Evaluate the following response for correctness. Consider whether it actually answers the question asked. Ignore presentation quality. Give a single number from 0 to 10, where 0 means entirely incorrect and 10 means fully correct. Output only that number.
\end{quote}
\emph{Referenced condition} (the default whenever a reference is shown):
\begin{quote}\small
(a) You are a strict grader. You are given a question, a reference answer known to be correct, and a candidate response. Rate the candidate's correctness on a scale from 0 (completely wrong) to 10 (fully correct), using the reference to check it. Judge the content, not the writing. Respond with only the integer score.\\[2pt]
(b) Compare the candidate response against the trusted reference answer provided. Score how correct the candidate is, from 0 (entirely incorrect) to 10 (fully correct). Presentation does not matter. Output only a single number.
\end{quote}
\emph{Mitigation conditions only.} The full-credit sentence and the second instruction wording each append one of these sentences to the referenced rubric:
\begin{quote}\small
Note: the reference shown is one valid answer among possibly many; the question may have several correct answers, and any correct answer deserves full credit even if it differs from the reference.\\[2pt]
Important: do not compare the candidate's answer to the reference. Use the reference only to understand the problem. Verify the candidate's answer on its own merits; a correct answer that differs from the reference must receive the same score as one that matches it.
\end{quote}
\emph{Blind-then-reveal condition.} The first call scores the candidate with the blind rubric. The second call appends to the referenced rubric: ``You have already assessed this candidate without seeing any reference. A correct answer that differs from the reference must keep the score you gave it; revise only if the reference reveals the candidate is actually wrong.'' Its user message shows the reference as ``Reference answer (one valid answer, trusted)'' and ends with ``Your blind score for this candidate was $x$ out of 10.''

\paragraph{What the default referenced rubrics ask.}
Both default referenced rubrics invite comparison with the reference: rubric~$a$ asks the judge to use the reference to check the candidate, and rubric~$b$ begins ``Compare the candidate response against the trusted reference answer provided.'' Neither says that a question can have several correct answers. The sentence ``the reference shown is one valid answer among possibly many'' is not part of either default rubric; it defines a separate mitigation condition, reported as such (\cref{app:acceptance,app:r24,app:mitigation}). A judge that follows the comparison instruction could therefore favor the reference because it was told to. The effect does not rest on our wording. It appears under the verbatim MT-Bench single-answer template (\cref{app:rotation}), which also asks for a comparison with the reference; under a direct prompt that asks only whether the response is correct and says nothing about how to use the reference (\cref{app:r26}); under chain-of-thought, verify-first, and solve-first prompts that ask the judge to use or check the reference, although reasoning shrinks it (\cref{app:r25}); for correct answers the judge had already accepted blind (\cref{app:acceptance}); and, under correctness prompts, for answers credited blind (\cref{app:r27}). The wording does change its size (\cref{app:r26}).

\paragraph{Other prompts.}
The benchmark graders are copied verbatim from their repositories (\cref{app:r24}); the direct, chain-of-thought, verify-first, and solve-first prompts are quoted in \cref{app:r25,app:r26}; the open RL-reward verifiers use their own templates (\cref{app:r26}).

\paragraph{The Figure 1 example.}
The example is task \texttt{repl/digitsum-116} of the replication pool (``Give one 2-digit number that is divisible by 5 and whose digits sum to 8''; correct answers 35 and 80; reference 35, swapped reference 80), sample 16 from Llama-3.1-8B-Instruct, scored by Qwen2.5-7B with the exact scorer. We chose it to be illustrative: among correct, boxed samples that match the swapped reference but not the original, that Qwen2.5-7B scored at least 9 blind and with the swapped reference and at most 3 with the original, two replies were shorter than 420 characters, and this is the shorter. Its swap (6.4 points) is larger than the pool average (\cref{fig:rotation-pool}). The full reply:
\begin{quote}\small
To find a 2-digit number that is divisible by 5, the ones digit must be either 0 or 5.\\[2pt]
Let's consider the case where the ones digit is 0. In this case, the tens digit must be 8 for the digits to sum to 8. The number would be 80.\\[2pt]
Now, let's check if 80 is divisible by 5. Since $80 \div 5 = 16$, 80 is indeed divisible by 5.\\[2pt]
The number 80 meets the given conditions, so we can stop here.\\[2pt]
The answer is \texttt{\textbackslash boxed\{80\}}.
\end{quote}

%==========================================================================================
\FloatBarrier
\section{How widely graders rely on one reference}
\label{app:survey}

We surveyed benchmarks whose official scorer is an LLM judge, evaluation libraries and cloud platforms with a built-in correctness judge, RL reward pipelines that use an LLM verifier, and research papers from 2023 to 2026 that evaluate or train with an LLM judge. Every entry was checked on 14 September 2026 against its primary source (repository code at a pinned commit, official documentation, or the paper); the records, 24 verbatim judge templates, and all quotes are in the supplementary material. An entry counts as \emph{reference-guided} when the judge sees one gold answer, or a closed list of them, and decides correctness from it (\cref{tab:survey}).

\begin{table}[htbp]
\centering
\caption{Surveyed graders that give an LLM judge one reference answer or a closed list of them. Research papers were selected because they use such a judge; 10 checked candidates that did not were dropped. Partly: Arize Phoenix (reference template only in an older release) and Rubrics as Rewards (references write the rubrics).\ph}
\label{tab:survey}
\footnotesize
\setlength{\tabcolsep}{3pt}
\begin{tabular}{L{2.2cm}ccL{7.8cm}}
\toprule
Group & Surveyed & One reference & Reference-guided entries \\
\midrule
Benchmarks & 16 & 11 & SimpleQA~\citep{wei2024simpleqa}, SimpleQA Verified~\citep{haas2025simpleqaverified}, BrowseComp~\citep{wei2025browsecomp}, HLE~\citep{phan2025hle}, the simple-evals MATH checker~\citep{simpleevals}, MT-Bench~\citep{zheng2023judging}, FreshQA~\citep{vu2023freshllms}, CRAG~\citep{yang2024crag}, MixEval~\citep{ni2024mixeval}, Omni-MATH~\citep{gao2024omnimath}, FRAMES~\citep{krishna2025frames}. Not: TruthfulQA, HealthBench, GAIA, LiveBench, Arena-Hard \\
Evaluation libraries & 12 & 11 + 1 partly & Inspect AI~\citep{inspectai}, OpenAI Evals~\citep{openaievals}, LangChain~\citep{langchainevals}, openevals~\citep{openevals}, LlamaIndex~\citep{llamaindexcorrectness}, RAGAS~\citep{es2023ragas}, DeepEval~\citep{deepeval}, MLflow~\citep{mlflowjudges}, promptfoo~\citep{promptfoofactuality}, TruLens~\citep{trulens}, OpenCompass~\citep{opencompassjudge}; Phoenix~\citep{phoenixevals} partly \\
Cloud platforms & 4 & 3 & OpenAI graders~\citep{openaigraders}, Azure AI Foundry~\citep{azurefoundrysimilarity}, AWS Bedrock~\citep{bedrockjudgeprompts}. Not: Vertex AI \\
RL reward pipelines & 11 & 10 + 1 partly & DeepSeek-R1~\citep{deepseek2025r1}, Kimi k1.5~\citep{kimi2025k15}, Qwen3~\citep{yang2025qwen3}, Seed1.5-Thinking~\citep{seed2025thinking}, General-Reasoner~\citep{ma2025generalreasoner}, RLVR across domains~\citep{su2025rewardbridge}, xVerify~\citep{chen2025xverify}, CompassVerifier~\citep{liu2025compassverifier}, TinyV~\citep{xu2025tinyv}, OpenAI reinforcement fine-tuning~\citep{openairft}; Rubrics as Rewards~\citep{gunjal2025rar} partly \\
Research papers, 2023--2026 & --- & 34 & Math and science~\citep{huang2025pitfalls,chernyshev2024umath,lu2025ineqmath,xu2025ugphysics,minimax2025m1,cheng2025guru,bercovich2025llamanemotron,huang2024olympicarena,fang2024mathodyssey}; proofs and competition grading~\citep{petrov2025proofbluff,luong2025imobench,mahdavi2025refgrader,balunovic2025matharena}; code~\citep{zhuo2024icescore,tong2024codejudge}; QA and RAG~\citep{chandak2025answermatching,he2024chinesesimpleqa,ru2024ragchecker,rau2024bergen,refguided2024,kamalloo2023evaluating,thakur2024judging,wang2023evouna}; agents and deep research~\citep{wu2025webwalker,zhou2025browsecompzh,wu2025webdancer,li2025websailor,gao2025asearcher,zheng2025deepresearcher,chen2025browsecompplus,du2025deepresearchbench,li2025webthinker,liu2025webexplorer,tongyi2025deepresearch} \\
\bottomrule
\end{tabular}
\end{table}

\paragraph{What the prompts say.} We could read 34 judge prompts in full (a different set from the 34 research papers). Of these, 27 define correctness as matching the reference, 4 partly, and 3 offer it only as an aid. Two say that several answers can be valid, and both give the judge a closed list~\citep{ni2024mixeval,vu2023freshllms}; three more allow different wording but require the reference's content. None tells the judge to credit a correct answer outside the reference. Five declare the reference correct: ``Take it as granted that the Ground Truth is always correct'' (CRAG), ``The reference answer is ALWAYS correct'' (Omni-MATH), ``THE STANDARD ANSWER IS ALWAYS CORRECT AND THE QUESTION IS PERFECTLY VALID. NEVER QUESTION THEM'' (CompassVerifier), ``You can assume that the standard answer is definitely correct'' (OpenCompass), and the HLE and BrowseComp judge's ``precise and unambiguous [correct\_answer]''. Reward prompts follow the same pattern: DeepSeek-R1's reward prompt rates an answer correct when it ``fully aligns with the reference answer''~\citep{deepseek2025r1}.

\paragraph{Research papers.} Read this group as an illustrative case series, not a census. Unlike the benchmarks, libraries, platforms and RL pipelines above, which are a bounded enumeration, the papers were gathered by area --- math and science, proof grading, code, QA and RAG, agents and deep research --- and screened to those that judge against a single reference or a closed list, with the dropped candidates recorded in the supplementary material. We fixed no query string, database or venue list in advance, so these papers show that the pattern recurs across subfields; they do not estimate how common it is in the literature. Eight of the 34 papers use the judge as an RL reward or data filter. In seven others, all evaluations, the tasks admit several valid answers and the judge still sees one reference or a closed list~\citep{kamalloo2023evaluating,wang2023evouna,zheng2025deepresearcher,du2025deepresearchbench,zhuo2024icescore,tong2024codejudge,petrov2025proofbluff}; among them, providing reference code made all three LLM-based evaluators worse on CoNaLa~\citep{tong2024codejudge}. IMO-ProofBench's grader is told the reference ``demonstrates one valid approach'' and is reported to be ``prone to be overly punitive'' for unconventional correct solutions~\citep{luong2025imobench}.

%==========================================================================================
\FloatBarrier
\section{How common are questions with several valid answers?}
\label{app:prevalence}

We sampled 200 questions (seed 0) from each of four datasets that ship one reference answer per question and labeled each as having multiple non-equivalent valid answers, a unique answer up to equivalence, or unclear. Equivalent forms of one answer (for example $1/2$ and $0.5$) count as unique, and ``find all'' questions whose answer is a single set count as unique. Each item was labeled twice independently by LLM annotators (Claude models) that did not see each other's labels, and disagreements were adjudicated by re-reading the item. Every alternative answer given for a math item was checked by computation. Both passes come from one model family, so agreement may overstate reliability, and we did not validate the labels against human annotators.
In this first audit (table in the extended appendix), curated math sets are almost entirely single-answer; the multi-answer math items are mostly underspecified problems rather than requests for an example. In NQ-open, 30.0\% of questions admit several correct answers, and 18.5\% after excluding questions whose answer changed over time.

\paragraph{Benchmarks graded against one reference, and a second annotator family.} A pre-specified extension (protocol committed before any labeling call) sampled 200 questions each from SimpleQA~\citep{wei2024simpleqa}, BrowseComp~\citep{wei2025browsecomp}, text-only exact-answer HLE questions~\citep{phan2025hle}, TriviaQA~\citep{joshi2017triviaqa}, and the WebInstruct-verified RL training set~\citep{ma2025generalreasoner}, with the definitions above. Two annotators from different developers labeled every item independently, Claude Sonnet~5 and GPT-5-mini; Claude Opus~5 adjudicated disagreements until its spending cap, after which the remaining 135 disagreements (106 on BrowseComp, most about whether an item needs deep solving) kept the conservative label (\cref{tab:prevalence-v2}); the fallback can only lower the count of questions with several answers, and had the six HLE fallbacks that involve that label gone the other way, HLE would reach 6.0\%. Questions written for a single answer rarely admit several: 0\% to 4\%. The RL training set reaches 5.5\%, and 18\% by either annotator. Agreement between the two annotators is low ($\kappa \le 0.35$ on the several-versus-other label), so these rates are rough; none was checked by humans. A GPT-5-mini relabel of the first audit's 800 items agrees with its final labels at $\kappa = 0.53$ on the several-versus-other label and marks more items as having several answers (87 of 200 NQ-open questions, against 60), so the first audit's rates are likely low.

\begin{table}[htbp]
\centering
\caption{Share of sampled questions with several valid answers after adjudication (95\% Wilson interval), agreement between the two annotators on the several-versus-other label, and the share labeled several by both and by either annotator. 200 questions per dataset; LLM annotators, not humans.\pre}
\label{tab:prevalence-v2}
\footnotesize
\begin{adjustbox}{max width=\linewidth}\begin{tabular}{lrrrrcccc}
\toprule
Dataset & $n$ & Several & Unique & Unclear & Several (\%) & $\kappa$ & Both (\%) & Either (\%) \\
\midrule
SimpleQA & 200 & 2 & 197 & 1 & 1.0 {\scriptsize$[0.3,3.6]$} & 0.00 & 0.0 & 1.5 \\
BrowseComp & 200 & 0 & 92 & 108 & 0.0 {\scriptsize$[0.0,1.9]$} & --- & 0.0 & 0.0 \\
HLE (text, exact answer) & 200 & 6 & 141 & 53 & 3.0 {\scriptsize$[1.4,6.4]$} & 0.10 & 1.0 & 15.0 \\
TriviaQA & 200 & 8 & 192 & 0 & 4.0 {\scriptsize$[2.0,7.7]$} & 0.29 & 1.5 & 8.0 \\
WebInstruct-verified (RL) & 200 & 11 & 173 & 16 & 5.5 {\scriptsize$[3.1,9.6]$} & 0.35 & 4.5 & 18.0 \\
\bottomrule
\end{tabular}
\end{adjustbox}
\end{table}

%==========================================================================================
\FloatBarrier
\section{Judge biases and optimization}
\label{app:relatedbias}

LLM judges favor the first response~\citep{wang2024fair}, longer responses~\citep{dubois2024lengthcontrolled}, their own outputs~\citep{panickssery2024llm}, and scores shown to them~\citep{kapetanovic2026anchoring}, where the anchor is a score rather than an answer; the concern carries over to rubric-based rewards, where in one surveyed pipeline the reference is used to write the rubric~\citep{gunjal2025rar}. Catalogues and benchmarks of judge reliability~\citep{koo2024benchmarking,ye2025justice,tan2025judgebench} do not vary which correct answer is shown. The effect resembles sycophancy toward a suggested answer~\citep{sharma2023sycophancy,turpin2023unfaithful}, except that the suggestion here is correct.

\paragraph{Single references and many correct answers.}
Reference-based metrics in translation reward outputs resembling the particular reference~\citep{freitag2020bleu,fomicheva2016reference}, and open-domain question answering under-credits correct answers missing from the gold list~\citep{min2020ambigqa,si2021whatsinaname,bulian2022tomayto,kamalloo2023evaluating}, which benchmarks and answer-equivalence methods target~\citep{amouyal2022qampari,li2024pedants}. Those benchmarks supply questions with many answers; the failure we measure is upstream of them, in the grader, and appears even when the list of correct answers is complete.

%==========================================================================================
\FloatBarrier
\section{Preregistration ledger}
\label{app:ledger}

Every confirmatory prediction was committed before its analysis was run, and for new model calls before the data existed; the verdict rule applied to them was adopted afterwards (below). Three experiments were registered on OSF with third-party timestamps\osf{} (links withheld for review; anonymized copies are in the supplementary material). The rest were committed to the project repository with a timestamp that predates the analysis\pre. Some of these rescored or re-analyzed outputs whose direction of effect was already known, such as the exact scorer (\cref{app:estimator}); the main text marks them with $\ddagger$ and this appendix labels them \pre. Analyses marked post hoc\ph{} carry no registration weight.

Verdicts follow one rule. \emph{Confirmed} means the 95\% interval excludes zero in the predicted direction, or a registered point prediction holds. \emph{Not confirmed} means the interval includes zero. \emph{Falsified} means the interval excludes zero in the opposite direction, or a point prediction fails. The rule was written into a registration made partway through the project for later experiments, and applied retroactively to all earlier predictions after most of their outcomes were known; several predictions that the registration documents mark as falsified are therefore reported here as not confirmed. The re-derivation moved predictions in both directions. Of these, 7 moved from falsified to not confirmed (P-R2B1 to P-R2B3, P-R3D1, P-R3D2, P-2D2, and the Qwen2.5-7B cell of P-RC1). Another 10 moved against us: P-M2 and P-S5 to falsified, P-FM3 from confirmed to falsified, P-S2, P-MC2, and P-EX2 to split, and P-FR1, P-FK2, P-SC1, and P-SC2 from confirmed to not evaluable.

A round is one dated batch of registered experiments; the registration documents in the supplement are numbered by round. Rounds 28 and 29 were drafted and never run --- an agentic-judge design and a human validation of our labels --- so they register no prediction and enter no total. Counting each registered prediction once, and from round 19 on once per judge or grader configuration, \cref{tab:ledger-totals} gives the totals. A prediction registered for several judges or training arms that got different verdicts across them is \emph{split}. \Cref{tab:ledger-failures} lists every prediction that was not confirmed, falsified, split, not evaluable, or not scored, and the excluded check; the extended appendix lists every registered experiment, including those whose predictions were all confirmed.

\begin{table}[htbp]
\centering
\caption{Registered predictions by verdict.}
\label{tab:ledger-totals}
\footnotesize
\begin{tabular}{lccccc}
\toprule
Verdict & To round 26 & R27 (\cref{app:r27}) & R30 (\cref{app:r30}) & R31 (\cref{app:r31}) & R32 (\cref{app:holdout}) \\
\midrule
Confirmed & 202 & 11 & 14 & 4 & 5 \\
Not confirmed & 14 & 0 & 0 & 4 & 0 \\
Falsified & 34 & 3 & 2 & 0 & 0 \\
Split & 17 & 0 & 1 & 0 & 0 \\
Not evaluable & 18 & 0 & 0 & 0 & 0 \\
Not scored & 2 & 0 & 0 & 0 & 0 \\
Excluded (\cref{app:frontier}) & 3 & 0 & 0 & 0 & 0 \\
\midrule
Total & 290 & 14 & 17 & 8 & 5 \\
\bottomrule
\end{tabular}
\end{table}

{\footnotesize
\setlength{\tabcolsep}{4pt}
\begin{longtable}{L{3.5cm}L{9.5cm}}
\caption{Every registered prediction that was not confirmed, falsified, split, not evaluable, or not scored, by experiment. Counts in parentheses where one entry covers several predictions. The rows contain 18 not confirmed, 39 falsified, 18 split, 18 not evaluable, and 2 not scored predictions, and the 3 excluded ones.}\label{tab:ledger-failures}\\
\toprule
Experiment & Predictions and outcome \\
\midrule
\endfirsthead
\multicolumn{2}{l}{\small\emph{\cref{tab:ledger-failures}, continued}} \\
\toprule
Experiment & Predictions and outcome \\
\midrule
\endhead
\bottomrule
\endfoot
$k$-pool and mitigations\osf{} & Blind-then-reveal keeps the $k{=}1$ benefit (\textbf{not confirmed}: its $k{=}1$ intervals include zero); validity gate removes the margin at every $k{>}1$ (\textbf{falsified}: Qwen2.5-7B $+5.5$ to $+6.7$; phi-4 $-0.51$ \ci{-1.01}{-0.12} at $k{=}4$). \\
Judge coverage\osf{} & Reference adds inflation for each added judge (\textbf{split}: gemma-2-9b and Qwen2.5-14B confirmed; Mistral-7B not confirmed, a score-resolution artifact); blind-then-reveal closer to blind than the reference (\textbf{split}: gemma-2-9b and Qwen2.5-14B confirmed; Mistral-7B falsified). \\
Resolution screen, Ministral-8B & Margin above Mistral-7B's (\textbf{falsified}: $+0.32$ against $+1.38$). \\
Synthetic grid & Margins rise with $k$ (\textbf{falsified}: flat, and falling for Qwen2.5-14B and phi-4); instruction lowers the margin at every $k$ (\textbf{split}: Qwen2.5-14B confirmed; Qwen2.5-1.5B, 3B, and 7B falsified); same pattern on a Llama-generated pool (\textbf{falsified}: no $k{=}1$ benefit). \\
Composition dose & Inflation rises with the number of matching valid answers and is null with none (\textbf{falsified}: it falls, and is largest with none). \\
Constructed lure test & Two predictions (\textbf{not scored}): closed without judging after its naturalness audit failed in five rounds. \\
Two-factor composition & Lures raise inflation when no answer matches (\textbf{split}: phi-4 confirmed, Qwen2.5-7B not confirmed); removing lures restores a rising dose response (\textbf{not confirmed}, both judges). \\
Distance gradient & Near, middle, and far thirds ordered (\textbf{falsified}, both judges). \\
Solution-route rotation (\cref{app:mechanism}) & No score change when only the derivation rotates (\textbf{split}: Qwen2.5-7B falsified, $-0.32$ \ci{-0.56}{-0.08}; phi-4 confirmed, $-0.02$ \ci{-0.04}{0.00}). \\
\textsc{MACE} (\cref{app:mace}) & Referenced margin positive with blind margin null (\textbf{split}: Qwen2.5-7B confirmed; phi-4 falsified, blind margin $-0.80$ \ci{-1.44}{-0.15}). \\
Counter-instructed swap (\cref{app:r31}) & Answers that stop matching are credited less under the second draw (\textbf{not confirmed}, both judges: $-0.007$ \ci{-0.031}{+0.015} for Sonnet~5 and $-0.021$ \ci{-0.051}{+0.007} for GPT-4.1). Answers that start matching are credited more (\textbf{not confirmed}, both judges: $+0.012$ \ci{-0.011}{+0.034} and $+0.000$ \ci{-0.016}{+0.018}). The swap does not survive an instruction not to use the reference. \\
Swap under the direct prompt (\cref{app:r30}) & Shift smaller than the same judge's shift under our rubric (\textbf{falsified} for GPT-4.1: $-0.264$ and $+0.247$ against $-0.125$ and $+0.189$; confirmed for Sonnet~5). Answers matching neither reference move at most 5 points (\textbf{falsified} for Sonnet~5: $-5.9$ \ci{-10.7}{-1.1}; confirmed for GPT-4.1). \\
Chosen reference (\cref{app:r30}) & Shortfall smaller than under the random draw (\textbf{split}: confirmed for the BrowseComp/HLE judge and our rubric in both arms, and for the SimpleQA grader under the modal rule; not confirmed for the SimpleQA grader under the first-listed rule, $+0.045$ \ci{-0.003}{+0.092}). \\
Frontier CLI judge (\cref{app:frontier}) & \textbf{Excluded} from the counts (3): the calls ran with the interface's default system prompt, tools, and repository context, and extended thinking engaged on 54\% of blind and 8\% of referenced calls. \\
Frontier API judges (\cref{app:frontier}) & Sonnet~5: \textsc{MACE} referenced margin positive with blind margin null (\textbf{falsified}: blind margin $-1.37$ \ci{-2.05}{-0.67}). GPT-4.1 failed the registered score-resolution screen, so its four inflation predictions are \textbf{not evaluable} (4); its $k$ step on the synthetic grid (\textbf{not confirmed}). \\
Code exemplar (extended appendix) & Canonical solution adds inflation where the test suite does not (\textbf{split}: Qwen2.5-7B not confirmed; Qwen2.5-14B falsified, it lowers inflation); similarity to the solution raises scores (\textbf{split}: Qwen2.5-7B confirmed; Qwen2.5-14B not confirmed). \\
Answer-format continuum (extended appendix) & Margin positive at integer, pair, and JSON formats; sharp drop at the program format; drop rather than gradient (\textbf{not evaluable} (3): 7 and 1 usable tasks at pair and JSON); inflation positive at the first three formats and not positive at the program format (\textbf{split}: phi-4 confirmed; Qwen2.5-7B falsified). \\
Reveal-call variants (\cref{app:mitigation}) & Neutral wording changes nothing (\textbf{split}: phi-4 confirmed, Qwen2.5-7B falsified). \\
Variants on gemma-2-9b (\cref{app:mitigation}) & Another judge's score works as well (\textbf{falsified}). \\
Training, runs 1--3 & Each arm raises its own judge's score: run 1 (\textbf{split}: reference arm confirmed, blind-then-reveal arm falsified); run 2 (\textbf{not confirmed}, both arms); run 3 (\textbf{split}: reference arm confirmed, blind-then-reveal arm falsified); run 3 blind arm (\textbf{falsified}). Reference arm copies the reference answer more (\textbf{falsified}: its rate fell). The other 10 predictions compare against an arm that failed its check (\textbf{not evaluable}, 10). \\
Decode cost (\cref{app:mechanism}) & All four predictions are \textbf{split} (4) because phi-4's cells are not evaluable (81\% to 98\% of its replies had no score). Qwen2.5-7B: margin survives one step (not confirmed); each further step lowers it (not confirmed); below half by three steps (confirmed); inflation non-increasing (falsified). \\
Constraint-problem pilot and follow-ups (extended appendix) & Selection harm, four designs: garbled invalid answers (\textbf{falsified}); no room to move, both conditions at zero (\textbf{not evaluable}); the reference's own text among the candidates (\textbf{falsified}); corrected design (\textbf{not confirmed}: right direction for both judges, intervals reach zero, 8 tasks); the corrected design pooled with a second valid solution per item, 16 comparisons (\textbf{split}: Qwen2.5-7B confirmed, $-0.31$ \ci{-0.56}{-0.13}; phi-4 not confirmed, $-0.13$ \ci{-0.31}{0.00}). \\
Blind-then-reveal on MATH-500 (\cref{app:mitigation}) & The protocol keeps less than half of the reference's benefit (\textbf{falsified}: it keeps nearly all of it). \\
Thinking judge (Qwen3-8B; \cref{app:rotation}) & Rotation gain, rotation loss, and referenced margin above blind at $k{>}1$ (\textbf{not confirmed}, 3); accuracy step below zero (\textbf{falsified}). No measurable anchoring. \\
\textsc{MACE} end to end (\cref{app:reasoning}) & Reference lowers accuracy at $k{>}1$ (\textbf{not confirmed}, both judges); blind-then-reveal beats the reference at $k{>}1$ (\textbf{not confirmed}, both judges); blind-then-reveal at least as accurate pooled over the natural mix (\textbf{falsified}, both judges). \\
Opus~5 on the Sonnet~5 subsamples (\cref{app:acceptance}) & Grading of correct non-matching answers worse with the reference on the replication pool (\textbf{not confirmed}); at most 5\% of wrong answers scored 10 blind in each pool (\textbf{falsified}: 5.1\% to 8.2\%). \\
Frontier-model label audit (\cref{app:grid}) & Opus~5 confirms at least 90\% of the replies the any-format rule credits (\textbf{falsified}, 86.2\%). \\
Benchmark graders, one reference and the full answer list (\cref{app:r24}) & Credit for wrong answers at most 3 points higher with the list (\textbf{falsified} (2) for the BrowseComp/HLE judge with GPT-4.1, $+3.6$, and with o3-mini, $+5.4$). \\
Thinking and the full-credit sentence together (\cref{app:r24}) & False rejections still rise (\textbf{not confirmed}: $-0.2$~pp \ci{-7.2}{+4.6}). \\
Judges that reason before deciding (\cref{app:r25}) & Referenced accuracy below the answer lists (\textbf{not confirmed} (2) for verify-first and solve-first GPT-4.1); credit for differing correct answers at least 10 points lower with the reference (\textbf{falsified}, 6 of 6: 0.8 to 8.6 points). \\
Instruction against reasoning (\cref{app:r26}) & Direct prompt blind within 5 points of the lists (\textbf{falsified} for GPT-4.1, $-5.4$); direct shortfall within 5 points of chain of thought (\textbf{falsified} for GPT-4.1, $-17.4$ against $-3.0$); removing the HLE judge's match clause leaves it within 10 points of the lists (\textbf{falsified}, both judges). \\
Grading with many wrong answers (\cref{app:r26}) & Grading AUROC with the reference no more than 0.05 below blind (\textbf{falsified} (4) for the direct and verify-first prompts, both judges). \\
Open RL-reward verifiers (\cref{app:r26}) & With the full list, within 5 points of the lists (\textbf{falsified} for General-Reasoner's verifier, $-24.2$). \\
Hard single-answer questions (\cref{app:r27}) & With the reference, solve-first and direct prompts credit at most 5 points more wrong answers than the verbatim HLE judge (\textbf{falsified} (2) for GPT-4.1, $+25.9$ and $+17.3$); solve-first recall at least the HLE judge's minus 5 points (\textbf{falsified} for Sonnet~5: 6 of 7 against 7 of 7). \\
\end{longtable}
}

%==========================================================================================
\FloatBarrier
\section{\textsc{MACE}: audits and matching}
\label{app:mace}

\textsc{MACE}~\citep{mace2026} asks factual questions whose number of correct answers is fixed by construction at $k \in \{1,2,4,6\}$. Before building the pool we audited each domain. All 2,000 Math-domain questions have correct and complete answer sets by direct computation, so the domain enters as a label-noise control. For the entity domains we measured how often a model names a valid answer missing from the shipped list: River (2.3\%) and Language (8.1\%) passed, and Region, Office, and Award (30.8\% to 43.5\%) were excluded.

\paragraph{Matching answers.} A candidate's answer is its last \texttt{\textbackslash boxed\{\}} content, or its last 200 characters if there is none, lowercased with non-alphanumeric characters removed. It is correct if it equals a normalized entry of the answer list, or contains an entry longer than three characters; the length filter gates containment alone, so a short name such as \emph{Po} or \emph{Wu} still matches exactly but never as a substring, which keeps it from matching inside an unrelated word. It matches the reference by the same test applied to the reference alone. A correct answer counted as differing from the reference therefore never contains the reference's name: an alias of the reference counts as a different correct answer only if it also contains another listed answer, and a reply that names several answers, the reference among them, counts as matching. Aliases missing from the list count as wrong, which lowers measured rejection of correct answers rather than raising it. Containment is also unanchored and we apply no negation check, so a reply that names an answer only to reject it counts as naming it.

\paragraph{Margins.} The step in added inflation from one to several correct answers holds in each audited domain, and the math control has the same sign.\pre{} Margins at $k{>}1$ on the audited domains are $+6.44$ for Qwen2.5-7B and $+5.86$ for phi-4, against blind margins of $-0.64$ and $-0.80$. Sonnet~5 has a referenced margin of $+5.29$ there, against a blind margin of $-1.37$. Per-domain tables and selection accuracy by $k$ are in the extended appendix.

%==========================================================================================
\FloatBarrier
\section{Reading answers: labelings and grading}
\label{app:grid}

\paragraph{Reading answers in any format.}
The registered checkers read an answer from its \texttt{\textbackslash boxed\{\}} content. On the $k$-pool a reply without a box is compared as a whole, and on the replication pool the checker needs exactly one integer in the reply's last 100 characters, so most answers stated in prose count as wrong. Only 61.6\% of replication samples and 57.1\% of $k$-pool samples use the box.
We relabeled every reply post hoc, before reading any judge score with the new labels. A reply counts as correct if the registered checker accepts it or the task's constraint accepts the extracted answer. In a second labeling, Qwen2.5-14B-Instruct reads each unboxed reply and reports the final answer it commits to; the rule and the model agree on 88\% of unboxed replies. A registered audit then had Claude Opus~5 (API, thinking disabled, the same extraction prompt) read a stratified random sample of 200 unboxed replies, 50 per pool and any-format label. Weighted by stratum sizes, its reading agrees with the any-format label on 91.1\% of unboxed replies (predicted at least 85\%; confirmed) and with the Qwen2.5-14B reading on 94.7\%. It confirms 86.2\% of the replies the any-format rule credits (predicted at least 90\%; falsified) and credits 7.4\% of those the rule rejects (predicted at most 15\%; confirmed). The any-format labels therefore over-credit some replies, which is one reason we also report boxed replies alone.

\paragraph{Grading under four labelings.}
\Cref{tab:labelings-gnm} reports grading AUROC for correct non-matching against wrong answers under the registered labels, both relabelings, and the registered labels restricted to boxed replies. Grading gets worse under the registered labels and both relabelings for every open judge; keeping only boxed replies shrinks the fall to 0.07 to 0.14 and removes it for Mistral-7B, whose integer scores failed our score-resolution screen. Under the registered labels and the exact scorer, Qwen2.5-7B scores 28.1\% of the replication pool's correct answers that differ from the reference 3 or below with the reference, against 0.4\% blind (93 tasks).\pre{} The selection loss at $k{>}1$ does not survive relabeling in the same way: crediting answers in any format leaves a loss that excludes zero only for Qwen2.5-7B, and with model-read labels no loss excludes zero. The preregistered inflation endpoints do survive relabeling: on the replication pool the pooled endpoint is $+1.19$ under the registered labels, $+1.14$ with answers read in any format, and $+1.12$ with model-read labels.\ph{} Selection-accuracy and false-rejection tables under all four labelings are in the extended appendix.

\begin{table}[htbp]
\centering
\caption{Change in per-task AUROC for separating correct answers that differ from the reference from wrong answers (95\% CI), under four labelings. Exact scorer for Qwen2.5 and gemma-2; matches are decided from the extracted answer. Qwen2.5-7B's replication value under the registered checker is pre-specified; the rest are post hoc.}
\label{tab:labelings-gnm}
\footnotesize
\setlength{\tabcolsep}{3pt}
\begin{adjustbox}{max width=\linewidth}\begin{tabular}{llcccc}
\toprule
Pool & Judge & Registered & Any format & Model-read & Boxed only \\
\midrule
Replication & Qwen2.5-7B & $-0.29$ {\scriptsize$[-0.32,-0.26]$} & $-0.27$ {\scriptsize$[-0.31,-0.24]$} & $-0.25$ {\scriptsize$[-0.28,-0.22]$} & $-0.10$ {\scriptsize$[-0.13,-0.08]$} \\
 & phi-4 & $-0.28$ {\scriptsize$[-0.31,-0.25]$} & $-0.24$ {\scriptsize$[-0.27,-0.20]$} & $-0.21$ {\scriptsize$[-0.24,-0.18]$} & $-0.08$ {\scriptsize$[-0.10,-0.06]$} \\
 & GPT-4.1 & $-0.04$ {\scriptsize$[-0.05,-0.02]$} & $-0.04$ {\scriptsize$[-0.06,-0.02]$} & $-0.03$ {\scriptsize$[-0.05,-0.02]$} & $-0.02$ {\scriptsize$[-0.04,-0.01]$} \\
\addlinespace
$k$-pool, $k{>}1$ & Qwen2.5-7B & $-0.30$ {\scriptsize$[-0.34,-0.25]$} & $-0.23$ {\scriptsize$[-0.28,-0.19]$} & $-0.21$ {\scriptsize$[-0.26,-0.16]$} & $-0.09$ {\scriptsize$[-0.15,-0.04]$} \\
 & Qwen2.5-14B & $-0.23$ {\scriptsize$[-0.28,-0.19]$} & $-0.17$ {\scriptsize$[-0.21,-0.14]$} & $-0.16$ {\scriptsize$[-0.20,-0.12]$} & $-0.07$ {\scriptsize$[-0.10,-0.04]$} \\
 & gemma-2-9b & $-0.23$ {\scriptsize$[-0.27,-0.19]$} & $-0.19$ {\scriptsize$[-0.23,-0.15]$} & $-0.16$ {\scriptsize$[-0.20,-0.12]$} & $-0.14$ {\scriptsize$[-0.19,-0.09]$} \\
 & phi-4 & $-0.33$ {\scriptsize$[-0.38,-0.27]$} & $-0.25$ {\scriptsize$[-0.29,-0.20]$} & $-0.22$ {\scriptsize$[-0.27,-0.17]$} & $-0.12$ {\scriptsize$[-0.18,-0.07]$} \\
 & Mistral-7B & $-0.10$ {\scriptsize$[-0.13,-0.06]$} & $-0.08$ {\scriptsize$[-0.11,-0.05]$} & $-0.08$ {\scriptsize$[-0.12,-0.04]$} & $+0.01$ {\scriptsize$[-0.05,+0.07]$} \\
 & GPT-4.1 & $-0.04$ {\scriptsize$[-0.08,-0.01]$} & $-0.01$ {\scriptsize$[-0.04,+0.00]$} & $-0.01$ {\scriptsize$[-0.03,+0.00]$} & $-0.01$ {\scriptsize$[-0.06,+0.02]$} \\
\bottomrule
\end{tabular}
\end{adjustbox}
\end{table}

%==========================================================================================
\FloatBarrier
\section{Exact scoring, a bare-answer reference, and selection}
\label{app:estimator}

\paragraph{Rescoring with the exact scorer.}
A pre-specified experiment rescored Qwen2.5-7B on the replication pool (blind and referenced), the swapped replication pool, the $k$-pool, \textsc{MACE}, and MATH-500, and Qwen2.5-14B, gemma-2-9b, and Mistral-7B on the $k$-pool, with the exact scorer of \cref{app:setup}. No Qwen2.5-7B reply fell below the captured-probability threshold, and the exact and original values differ by 0.07 to 0.17 points per score on average. The swap effect is slightly larger under the exact scorer and larger still on greedy integers (extended appendix), so the original rule did not create it.\pre
\paragraph{A bare-answer reference on MATH-500.}
The same experiment added a MATH-500 condition that shows only the final answer, with the rubric's ``reference solution'' changed to ``reference answer'' (\cref{tab:r4-math}).\pre

\begin{table}[htbp]
\centering
\caption{MATH-500, the 105 problems whose pool pass rate lies in $[0.2,0.8]$: best-of-64 selection accuracy (\%) blind, with the reference solution, and with its final answer alone, and the differences (points, 95\% CI). Exact scorer; both judges' scores are new runs, so blind accuracies differ slightly from the original-scorer files (68.6\% and 66.7\% against 67.6\%).\pre}
\label{tab:r4-math}
\footnotesize
\setlength{\tabcolsep}{3pt}
\begin{adjustbox}{max width=\linewidth}\begin{tabular}{lcccccc}
\toprule
Judge & Blind & Solution & Answer & Solution $-$ blind & Answer $-$ blind & Answer $-$ solution \\
\midrule
Qwen2.5-7B & 68.6 & 99.1 & 97.1 & $+30.5$ {\scriptsize$[+21.9,+40.0]$} & $+28.6$ {\scriptsize$[+18.1,+38.1]$} & $-1.9$ {\scriptsize$[-4.8,+0.0]$} \\
Qwen2.5-14B & 66.7 & 95.2 & 97.3 & $+28.6$ {\scriptsize$[+18.1,+38.1]$} & $+30.6$ {\scriptsize$[+22.1,+39.2]$} & $+2.1$ {\scriptsize$[-1.8,+6.0]$} \\
\bottomrule
\end{tabular}
\end{adjustbox}
\end{table}

\paragraph{Selection is less exposed than grading.}
On \textsc{MACE}'s multi-answer questions, best-of-64 accuracy did not fall ($+2.0$ and $+6.1$~pp for Qwen2.5-7B and phi-4, intervals including zero), so our registered loss was not confirmed (\cref{app:reasoning,app:mitigation}). The two endpoints measure different things, which is why they can disagree: grading scores every answer, while best-of-$n$ reports only the winner. One candidate that happens to match the drawn reference is lifted by that match, wins the pool, and is itself correct, so top-1 accuracy survives intact while every other correct answer in the same pool still carries the penalty. A pool wide enough to contain such a candidate therefore buffers selection against an effect that grading measures directly, and the losses concentrate, as above, where no candidate matches. On the $k$-pool at $k{>}1$, the reference lowers accuracy under the registered labels and the exact scorer for Qwen2.5-7B, Qwen2.5-14B, gemma-2-9b, and Mistral-7B (a re-analysis registered after the original-scorer accuracy readout had been seen), but with answers read in any format the loss excludes zero only for Qwen2.5-7B (\cref{app:grid}).

\paragraph{Blind verification.}
Blind verification is the AUROC over all answers without a reference. We predicted that it would be lower on \textsc{MACE} than on both constructed pools, that grading of correct non-matching answers would be worse on \textsc{MACE} and on the replication pool, and that overall grading would be better on \textsc{MACE}; all four were confirmed.\pre{} Exact-scorer and post hoc numbers quoted in the main text that appear in no table here are collected in the extended appendix.

%==========================================================================================
\FloatBarrier
\section{Reasoning, verification, and thinking}
\label{app:reasoning}
\paragraph{Thinking on versus off.}
This pre-specified experiment reran the Qwen3-8B protocol of \cref{app:rotation} (rubric~$a$ ending in a ``Score: N'' line, temperature 0.6, up to 4,096 new tokens) with the chat template's thinking switch off and nothing else changed, on the swap subsample, a 16-sample $k$-subsample, and an 81-task \textsc{MACE} subsample (table in the extended appendix). With thinking off, the loss under the swap is $-1.98$; with it on, $-0.12$. The \textsc{MACE} margin increase is $+7.07$ off and $+3.13$ on; a second toggle with Qwen3-14B repeats the pattern. With thinking on, valid swap samples receive 10 in 99.8\% of referenced replies and 96.2\% of blind ones, which leaves no room for a reference to raise them; a ceiling does not explain the missing loss, however, because the same answers keep their scores when the reference moves off them ($+0.01$). On \textsc{MACE}, away from the ceiling, the anchoring returns.

\paragraph{Sonnet~5 with and without thinking.}
A pre-specified three-arm design ran Sonnet~5 through the Claude command-line tool, with the default system prompt replaced, no tools, and an empty working directory, on 60 replication-pool tasks. The score-only arm anchors ($+3.57$ and $-2.85$ for text the reference moves onto and off; 26.7\% false rejections). Asking it to reason before a final ``Score: N'' line removes most of this on these arithmetic tasks, with thinking off ($+0.28$, $-0.20$, 0.4\% false rejections) or on ($+0.18$, $-0.32$, 0.4\%).\pre{} On 40 factual \textsc{MACE} tasks, the added margin is $+3.98$ \ci{+2.70}{+5.30} with a score only and $+3.03$ \ci{+1.70}{+4.65} with thinking on, and the share of non-matching correct answers scored 3 or below rises by 17.5 and 12.5 points. Blind, Sonnet~5 ranks both correct answers above the wrong one 94.4\% and 96.9\% of the time, so on these questions it favors the reference although it separates correct from wrong answers without one.\pre{} These command-line figures are superseded: the run stopped at the account's spend limit, and the complete pre-specified rerun through the Messages API gives 35.6 and 14.9 points on the same 40 tasks (\cref{app:acceptance}). Read the API numbers as the registered result and these as the partial readout they replaced.

\paragraph{Does acceptance mean verification?}
Keeping answers a judge scored 7 or above without a reference is a weak test of verification for the open judges: that filter keeps 93\% of correct replication-pool answers for Qwen2.5-7B but also 24\% of wrong ones, and 96\% and 30\% for phi-4 (any-format labels). It is closer to verification for Sonnet~5, which scores 7 or above on 7\% to 16\% of wrong answers in its subsamples, and for GPT-4.1 (12\%). Across the 120 replication tasks, Qwen2.5-7B's swap size falls as its blind AUROC rises (Spearman $-0.40$ \ci{-0.53}{-0.25}): it is 5.2 and 5.4 points in the lower two thirds of tasks (mean blind AUROC 0.83 and 0.93) and 3.6 points in the top third (0.98). Verification is associated with a smaller effect but, in these data, not with its absence.\ph

\paragraph{End to end on \textsc{MACE}.}
A pre-specified experiment measured best-of-64 selection accuracy on the benchmark's natural mix of $k$ (133 River and Language tasks). We predicted that the reference would lower accuracy at $k{>}1$, that blind-then-reveal would beat it there, and that blind-then-reveal would be at least as accurate over the natural mix; all three predictions failed for both judges.

%==========================================================================================
\FloatBarrier
\section{Swap test}
\label{app:rotation}

The ledger calls the swap test \emph{rot\-ation}: the reference is moved to a different correct answer while the candidate text stays fixed. The swap assigns each task the valid alternative that appears most often among its valid samples, so that enough samples match the new reference; the paired same-sample changes in \cref{tab:rotation} compare each sample with itself and are not affected by that choice. An earlier run on an exploratory pool built from \textsc{RewardBench-2} Ties found the effect largest for small Qwen2.5 judges and smaller but present at 14B (extended appendix).

\begin{table}[htbp]
\centering
\caption{Change in score of identical text when the reference is swapped (replication pool, 0--10 scale, 95\% CI; original scorer, registered runs).\pre{} \Cref{fig:rotation-pool} quotes the exact scorer, under which Qwen2.5-7B's ``Neither'' cell is $+0.12$, with an interval that includes zero; the full exact-scorer table is in the extended appendix. ``Onto'': samples that match the new reference. ``Off'': samples that matched only the old one. ``Neither'': valid samples matching neither. The last row keeps only samples Sonnet~5 scored at least 7 blind (3,194 of 3,494). For the open judges the analogous filter on each judge's own blind score gives $+4.99$ \ci{+4.69}{+5.32} (Qwen2.5-7B, 91\% kept) and $+4.08$ (phi-4, 98\% kept).\ph}
\label{tab:rotation}
\footnotesize
\setlength{\tabcolsep}{3pt}
\begin{adjustbox}{max width=\linewidth}\begin{tabular}{llcccc}
\toprule
Judge & Rubric & Moved onto ($n{=}120$) & Moved off ($n{=}93$) & Neither ($n{=}113$) & Onto minus neither \\
\midrule
Qwen2.5-7B & $a$ & $+4.97$ {\scriptsize$[+4.66,+5.27]$} & $-5.04$ {\scriptsize$[-5.38,-4.67]$} & $+0.13$ {\scriptsize$[-0.04,+0.31]$} & $+4.81$ {\scriptsize$[+4.44,+5.17]$} \\
 & $b$ & $+6.34$ {\scriptsize$[+6.10,+6.60]$} & $-6.58$ {\scriptsize$[-6.89,-6.26]$} & $+0.10$ {\scriptsize$[-0.11,+0.29]$} & $+6.18$ {\scriptsize$[+5.84,+6.50]$} \\
phi-4 & $a$ & $+4.09$ {\scriptsize$[+3.79,+4.38]$} & $-3.92$ {\scriptsize$[-4.26,-3.59]$} & $+0.30$ {\scriptsize$[+0.04,+0.57]$} & $+3.88$ {\scriptsize$[+3.52,+4.25]$} \\
 & $b$ & $+3.25$ {\scriptsize$[+2.95,+3.51]$} & $-3.18$ {\scriptsize$[-3.46,-2.87]$} & $+0.24$ {\scriptsize$[+0.01,+0.47]$} & $+3.08$ {\scriptsize$[+2.76,+3.41]$} \\
Claude Sonnet 5 & $a$ & $+3.04$ {\scriptsize$[+2.53,+3.56]$} & $-2.66$ {\scriptsize$[-3.36,-2.00]$} & $+0.04$ {\scriptsize$[-0.41,+0.46]$} & $+3.08$ {\scriptsize$[+2.45,+3.78]$} \\
 & $b$ & $+2.54$ {\scriptsize$[+2.04,+3.08]$} & $-2.51$ {\scriptsize$[-3.17,-1.89]$} & $+0.28$ {\scriptsize$[-0.10,+0.69]$} & $+2.31$ {\scriptsize$[+1.75,+2.86]$} \\
\quad blind-accepted only & $a$ & $+2.79$ {\scriptsize$[+2.26,+3.40]$} & $-2.50$ {\scriptsize$[-3.20,-1.83]$} & $+0.20$ {\scriptsize$[-0.21,+0.65]$} & $+2.67$ {\scriptsize$[+2.04,+3.33]$} \\
\bottomrule
\end{tabular}
\end{adjustbox}
\end{table}

\paragraph{Judges that explain or think before scoring.}
A pre-specified experiment re-ran the swap on a subsample of the replication pool (up to three samples per class per task, 879 samples) under two protocols: the verbatim MT-Bench single-answer templates with Qwen2.5-7B (the reference version asks the judge to begin ``by comparing the assistant's answer with the reference answer'' and to explain before a rating from 1 to 10), and Qwen3-8B with thinking enabled (\cref{tab:r2b}). Replies without a parseable score were dropped (20 of 879 swap pairs for MT-Bench). The thinking judge gives a score of 10 to 96.2\% of valid swap samples blind and to 99.8\% with either reference; under the MT-Bench template, by contrast, the mean score of valid swap samples falls from 6.68 blind to 4.97 with the reference.\ph

\begin{table}[htbp]
\centering
\caption{Judges that write before scoring: change in score of identical text under the swap (points, 95\% CI).\pre}
\label{tab:r2b}
\footnotesize
\setlength{\tabcolsep}{3pt}
\begin{adjustbox}{max width=\linewidth}\begin{tabular}{lccc}
\toprule
Protocol & Moved onto & No longer matches & Neither \\
\midrule
MT-Bench template, Qwen2.5-7B & $+3.618$ {\scriptsize$[+3.401,+3.842]$} & $-3.591$ {\scriptsize$[-3.936,-3.242]$} & $+0.077$ {\scriptsize$[-0.212,+0.404]$} \\
Thinking judge, Qwen3-8B & $-0.006$ {\scriptsize$[-0.017,+0.000]$} & $+0.011$ {\scriptsize$[+0.000,+0.032]$} & $-0.003$ {\scriptsize$[-0.009,+0.000]$} \\
\bottomrule
\end{tabular}
\end{adjustbox}
\end{table}

%==========================================================================================
\FloatBarrier
\section{Mechanism details}
\label{app:mechanism}

\paragraph{How each reference form is written.} The four forms change only the reference string the judge sees, never the candidate. \emph{Numeral} is the reference as stored. \emph{Words} renders every integer in it in English (``31'' becomes ``thirty-one''; lists element-wise, joined with commas), which is the same answer with no numeral in common with any candidate. \emph{Expression} renders every integer $n$ as $(n{-}1){+}1$ (``31'' becomes ``30+1''), so the answer is unchanged but equality now costs one computation step. For the worked-solution arm the reference is a full solution and \emph{masked} hides its result: the boxed content becomes \texttt{hidden}, and in the derivation every word-boundary occurrence of the answer is replaced. The worked columns use the 21 $k$-pool tasks whose answer is a single number because masking is only airtight there; for list-valued answers the derivation can still name the elements one by one, so the registered prediction was scored on the scalar families alone rather than on a leaky condition.

\begin{table}[htbp]
\centering
\caption{Masking suite: margin by how the reference is written (0--10 scale, 95\% CI). Numeral, words, and expression use 93 replication-pool tasks; the worked-solution columns use 21 $k$-pool tasks with a single-number answer.\pre}
\label{tab:masking}
\footnotesize
\setlength{\tabcolsep}{3pt}
\begin{adjustbox}{max width=\linewidth}\begin{tabular}{llcccc}
\toprule
Judge & Rubric & Numeral & Words & Expression & Worked: visible $\to$ masked \\
\midrule
Qwen2.5-7B & $a$ & $+4.62$ {\scriptsize$[+4.32,+4.96]$} & $+2.88$ {\scriptsize$[+2.47,+3.28]$} & $+0.29$ {\scriptsize$[+0.07,+0.49]$} & $+2.97$ {\scriptsize$[+2.01,+4.06]$} $\to$ $+0.05$ {\scriptsize$[-0.52,+0.63]$} \\
 & $b$ & $+5.99$ {\scriptsize$[+5.69,+6.29]$} & $+4.60$ {\scriptsize$[+4.14,+5.03]$} & $+0.35$ {\scriptsize$[+0.12,+0.56]$} & $+4.62$ {\scriptsize$[+3.57,+5.78]$} $\to$ $+0.31$ {\scriptsize$[-0.46,+1.13]$} \\
phi-4 & $a$ & $+3.79$ {\scriptsize$[+3.48,+4.11]$} & $+3.68$ {\scriptsize$[+3.27,+4.06]$} & $+0.89$ {\scriptsize$[+0.64,+1.13]$} & $+3.51$ {\scriptsize$[+3.06,+4.01]$} $\to$ $-0.49$ {\scriptsize$[-1.27,+0.26]$} \\
 & $b$ & $+3.11$ {\scriptsize$[+2.83,+3.39]$} & $+3.01$ {\scriptsize$[+2.71,+3.37]$} & $+0.63$ {\scriptsize$[+0.45,+0.82]$} & $+2.84$ {\scriptsize$[+2.40,+3.28]$} $\to$ $-0.45$ {\scriptsize$[-1.11,+0.12]$} \\
\bottomrule
\end{tabular}
\end{adjustbox}
\end{table}

\paragraph{Decode cost.}
We kept the answer an integer and re-encoded the reference with one, two, or three reversible steps, stating the decoding rule each time. We had predicted a graded decline; for Qwen2.5-7B the margin instead collapsed at the first step and stayed at zero. Every encoded prompt announced the encoding, which may explain the collapse by itself. The phi-4 cells are not evaluable: 81.0\% to 97.6\% of its encoded-reference replies contained no score.

\paragraph{Solution route.}
Holding the reference answer fixed and swapping its derivation to a different method moves the identical sample's score by $-0.32$ and $-0.40$ points for Qwen2.5-7B and $-0.02$ for phi-4; the effects are 20 to 400 times smaller than swapping the answer.

%==========================================================================================
\FloatBarrier
\section{Frontier judges}
\label{app:frontier}

\paragraph{API judges.}
Sonnet~5 reproduces the swap (\cref{tab:rotation}), the blind-accepted filter, the masking pattern (bare reference $+3.06$, words $+3.14$, expression $+0.47$), and the \textsc{MACE} result that its blind margin is negative. GPT-4.1 picks matching answers far more often with the reference (32 referenced-only against 1 blind-only picks) and has a margin of $+1.19$. It failed our pre-specified score-resolution screen: an average of 15.03 candidates per pool tie at its maximum score, against a bar of 3.0, so we report its inflation as unresolved. An earlier command-line check (August 2026) ran with the interface's default agent context, which may have shown repository metadata naming the hypothesis; it is excluded from the ledger counts and the main text and described in the extended appendix.

\paragraph{Sonnet~5 with wrong answers.}
A pre-specified experiment scored 16-candidate subsamples that include wrong answers, blind and with the reference: 40 replication tasks (8 per family), 40 \textsc{MACE} tasks with $k{>}1$, and all 22 eligible $k$-pool tasks with $k{>}1$, each subsample holding at least one matching-correct, one non-matching-correct, and one wrong answer. It began through the command-line interface and stopped at the account's spend limit; a registered amendment re-ran every call through the Messages API with extended thinking disabled and the same prompts, and those scores give the results (\cref{fig:onevsseveral}; table in the extended appendix).\pre{} Before the amendment we had seen the partial command-line readouts on 24 replication tasks, which the amendment records. The 40 \textsc{MACE} subsamples are the ``40-task subsample'' of the main text.

%==========================================================================================
\FloatBarrier
\section{Acceptance, the full-credit sentence, and benchmark scores}
\label{app:acceptance}

\paragraph{Opus~5 on the Sonnet~5 subsamples.}
A pre-specified experiment ran Claude Opus~5 through the Messages API with thinking disabled, with the prompts and subsamples of \cref{app:frontier}. Grading AUROC for correct non-matching answers changes by $-0.164$ \ci{-0.223}{-0.110} on \textsc{MACE}, $-0.057$ on the $k$-pool, and $-0.015$ \ci{-0.038}{+0.009} on the replication pool, and false rejections on \textsc{MACE} rise by 36.3 points \ci{+26.7}{+46.7}. Pooled over the three pools, 76 of 525 non-matching correct answers scored 10 blind fall to 3 or below with the reference, against 0 of 247 matching ones. We predicted that at most 5\% of wrong answers would score 10 blind in each pool; 5.9\%, 5.1\%, and 8.2\% did, so a blind 10 is a slightly weaker sign of verification for Opus~5 than for Sonnet~5 (1.8\% to 3.7\%).

\paragraph{Sonnet~5 with thinking through the API.}
A pre-specified experiment re-ran the 40 \textsc{MACE} subsamples with adaptive thinking (1,280 calls, none failed or unparseable). Grading AUROC changes by $-0.091$ \ci{-0.140}{-0.050}, against $-0.153$ with thinking disabled, and false rejections rise by 14.9 points \ci{+8.2}{+23.1}, against 35.6; all three predictions were confirmed.

\paragraph{A second blind call.}
A pre-specified experiment re-scored blind, through the same interface and with the same settings, every correct answer that each API judge had scored 10 on its first blind call (\cref{tab:r20-retest}). On the second call Sonnet~5 gives 3 or below to none of the 318 non-matching answers and still gives 10 to 87.1\% of them; Opus~5 gives 3 or below to 2 of 525 and 10 to 96.2\%. All three predictions were confirmed.

\begin{table}[htbp]
\centering
\caption{Correct answers each API judge scored 10 on its first blind call: how many score 3 or below with the reference, and on a second blind call through the same interface. One of Opus~5's second calls, on a matching \textsc{MACE} answer, returned no score, so its pooled matching count is 246 here and 247 in the registered acceptance readout.\pre}
\label{tab:r20-retest}
\footnotesize
\begin{adjustbox}{max width=\linewidth}\begin{tabular}{llrrrrrr}
\toprule
 & & \multicolumn{3}{c}{Differ from the reference} & \multicolumn{3}{c}{Match the reference} \\
\cmidrule(lr){3-5}\cmidrule(lr){6-8}
Judge & Pool & Blind 10 & Ref.\ $\leq 3$ & Retest $\leq 3$ & Blind 10 & Ref.\ $\leq 3$ & Retest $\leq 3$ \\
\midrule
Sonnet 5 & \textsc{MACE}, $k{>}1$ & 81 & 15 (18.5\%) & 0 (0.0\%) & 38 & 0 (0.0\%) & 0 (0.0\%) \\
 & Replication & 150 & 13 (8.7\%) & 0 (0.0\%) & 60 & 0 (0.0\%) & 0 (0.0\%) \\
 & $k$-pool, $k{>}1$ & 87 & 13 (14.9\%) & 0 (0.0\%) & 39 & 0 (0.0\%) & 0 (0.0\%) \\
 & Pooled & 318 & 41 (12.9\%) & 0 (0.0\%) & 137 & 0 (0.0\%) & 0 (0.0\%) \\
\addlinespace
Opus 5 & \textsc{MACE}, $k{>}1$ & 201 & 63 (31.3\%) & 2 (1.0\%) & 102 & 0 (0.0\%) & 0 (0.0\%) \\
 & Replication & 211 & 10 (4.7\%) & 0 (0.0\%) & 88 & 0 (0.0\%) & 0 (0.0\%) \\
 & $k$-pool, $k{>}1$ & 113 & 3 (2.7\%) & 0 (0.0\%) & 56 & 0 (0.0\%) & 0 (0.0\%) \\
 & Pooled & 525 & 76 (14.5\%) & 2 (0.4\%) & 246 & 0 (0.0\%) & 0 (0.0\%) \\
\addlinespace
Both judges & Pooled & 843 & 117 (13.9\%) & 2 (0.2\%) & 383 & 0 (0.0\%) & 0 (0.0\%) \\
Both judges & \textsc{MACE}, $k{>}1$ & 282 & 78 (27.7\%) & 2 (0.7\%) & 140 & 0 (0.0\%) & 0 (0.0\%) \\
\bottomrule
\end{tabular}
\end{adjustbox}
\end{table}

\paragraph{A matched acceptance threshold for the open judges.}
Open judges' expected scores are deterministic, so a second call would repeat the first, and they rarely give any answer exactly 10. For each open judge and pool we took the lowest blind expected score, on a grid of 0.05, that at most 3\% of wrong answers reach, and asked how often correct answers at or above it score 3 or below with the reference (\cref{fig:reasoning}; table in the extended appendix). On the constructed pools, 5\% to 57\% (by judge and pool) of those that differ from the reference score 3 or below, against at most 0.4\% of those that match; on \textsc{MACE}, 83.3\% (Qwen2.5-7B) and 41.9\% (phi-4), against none of the matching ones.\ph{} For Qwen3 on the \textsc{MACE} subsample we use a blind 10, which with thinking on 21\% to 24\% of wrong answers also receive, so those rows do not isolate acceptance.
\paragraph{How the two thresholds were chosen, and what moving them does.}
Two cut-offs recur in this paper, and neither is derived from theory. An answer counts as accepted when it reaches the blind score that at most 3\% of wrong answers reach, and as rejected when it scores 3 or below out of 10. We chose both for interpretability, so we report what happens when they move rather than defending the particular values. The acceptance cut-off is calibrated separately for each judge and pool, which is what makes it comparable across judges whose score distributions differ; a fixed cut would mean something different for each judge. The paragraph above gives the resulting rejection shares, and the binned analysis in the extended appendix shows the whole gradient behind them: the share of differing correct answers rejected is lower in a judge's top blind-score bin than in its lowest in all fourteen judge--pool cells, although the decline is monotone in only four of them, and it stays above zero in every top bin, while matching answers are rejected at more than 3\% only in each judge's lowest bin, except in bins of fewer than 15 answers. Moving the acceptance cut-off barely changes what it shows. Letting 5\% of wrong answers through instead of 3\% raises the \textsc{MACE} rejection shares from 83.3\% to 85.0\% for Qwen2.5-7B and from 41.9\% to 43.0\% for phi-4.\ph{} On the rejection side, \cref{tab:r21-graders} reports the benchmark study at score thresholds 5, 7, and 10 rather than at a single cut. The reference falls short of the answer lists at all three. The comparison with blind judging does move, however: at a threshold of 10 the reference beats blind judging for Sonnet~5 and Haiku~4.5, though not for Opus~5, because a blind judge rarely awards a 10. The shortfall against the lists is what we claim; the margin over blind judging depends on where the cut-off sits.

\paragraph{The full-credit sentence, measured by grading.}
A pre-specified experiment appended the full-credit sentence of \cref{app:setup} to the referenced rubric and re-ran the referenced calls on the Sonnet~5 \textsc{MACE} and replication subsamples and the Opus~5 \textsc{MACE} subsample, reusing each judge's blind calls. All four predictions were confirmed, those on the replication pool by point estimates alone: with the sentence, grading on \textsc{MACE} stays below blind for both judges, and Sonnet~5's rise in false rejections is smaller than without the sentence on both pools. Against the plain reference, the sentence raises \textsc{MACE} grading AUROC by 0.09 (Sonnet~5) and 0.11 (Opus~5) and lowers false rejections by 17.6 and 28.0 points; on the replication pool neither change is distinguishable from zero.\ph

\paragraph{Benchmark scores with our rubric.}
A pre-specified experiment asked Claude Haiku~4.5, Sonnet~5, and Opus~5 for one answer to each of the 150 River and Language questions with several correct answers, and took Llama-3.1-8B-Instruct's first stored sample for each; Opus~5 returned no text on one question, which leaves 149. Answers were labeled with the \textsc{MACE} answer lists and matching rule. Sonnet~5, Opus~5, and Haiku~4.5 (API, thinking disabled) scored each blind and with the reference, a score of 7 or above counting as correct (\cref{tab:r21-graders}); all 12 predictions were confirmed. With the reference, pooled accuracy falls below the answer lists under every grader; blind, Opus~5 differs from the lists by $+0.4$ points and Haiku~4.5, a weaker grader, by $-9.2$. Opus~5 left 11 referenced replies without a score, cut off at 512 tokens while it reasoned about whether a differing answer was acceptable; their questions drop out of its readouts (141 remain), which may understate its shortfall. Sonnet~5 grades its own answers here; the other three models' answers are free of that concern, and all four show the shortfall. GPT-4.1 with our rubric is in \cref{app:r24}.

\begin{table}[htbp]
\centering
\caption{The benchmark study under three graders: pooled changes against the answer lists (points, 95\% CI) at score thresholds 5, 7, and 10, the max rule at 7, and the shares of correct non-matching answers scored 3 or below, blind and with the reference. For Opus~5 and Haiku~4.5, threshold 7 and the max rule carry registered predictions and thresholds 5 and 10 are registered readouts; Sonnet~5's thresholds 5 and 10 and its max rule are post hoc.}
\label{tab:r21-graders}
\footnotesize
\begin{tabular}{lrcccc}
\toprule
Grader & Threshold & Blind $-$ list & Ref.\ $-$ list & Max $-$ list & $\leq 3$: blind; ref. \\
\midrule
Sonnet 5 & 5 & \begin{tabular}[c]{@{}c@{}}$+0.2$\\{\scriptsize$[-1.3,+1.7]$}\end{tabular} & \begin{tabular}[c]{@{}c@{}}$-17.6$\\{\scriptsize$[-21.8,-13.3]$}\end{tabular} & --- & 1.0\%; 24.2\% \\
 & 7 & \begin{tabular}[c]{@{}c@{}}$-1.7$\\{\scriptsize$[-3.5,+0.0]$}\end{tabular} & \begin{tabular}[c]{@{}c@{}}$-22.3$\\{\scriptsize$[-27.0,-17.9]$}\end{tabular} & \begin{tabular}[c]{@{}c@{}}$-1.0$\\{\scriptsize$[-2.5,+0.5]$}\end{tabular} & 1.0\%; 24.2\% \\
 & 10 & \begin{tabular}[c]{@{}c@{}}$-53.2$\\{\scriptsize$[-58.4,-47.8]$}\end{tabular} & \begin{tabular}[c]{@{}c@{}}$-43.5$\\{\scriptsize$[-49.0,-37.9]$}\end{tabular} & --- & 1.0\%; 24.2\% \\
\addlinespace
Opus 5 & 5 & \begin{tabular}[c]{@{}c@{}}$+0.5$\\{\scriptsize$[-0.5,+1.8]$}\end{tabular} & \begin{tabular}[c]{@{}c@{}}$-14.9$\\{\scriptsize$[-18.8,-11.0]$}\end{tabular} & --- & 0.5\%; 23.4\% \\
 & 7 & \begin{tabular}[c]{@{}c@{}}$+0.4$\\{\scriptsize$[-0.7,+1.6]$}\end{tabular} & \begin{tabular}[c]{@{}c@{}}$-15.6$\\{\scriptsize$[-19.9,-11.5]$}\end{tabular} & \begin{tabular}[c]{@{}c@{}}$+0.5$\\{\scriptsize$[-0.5,+1.8]$}\end{tabular} & 0.5\%; 23.4\% \\
 & 10 & \begin{tabular}[c]{@{}c@{}}$-4.8$\\{\scriptsize$[-7.1,-2.7]$}\end{tabular} & \begin{tabular}[c]{@{}c@{}}$-22.7$\\{\scriptsize$[-27.1,-18.4]$}\end{tabular} & --- & 0.5\%; 23.4\% \\
\addlinespace
Haiku 4.5 & 5 & \begin{tabular}[c]{@{}c@{}}$-6.5$\\{\scriptsize$[-9.6,-3.7]$}\end{tabular} & \begin{tabular}[c]{@{}c@{}}$-14.4$\\{\scriptsize$[-18.6,-10.2]$}\end{tabular} & --- & 9.1\%; 25.2\% \\
 & 7 & \begin{tabular}[c]{@{}c@{}}$-9.2$\\{\scriptsize$[-12.6,-6.0]$}\end{tabular} & \begin{tabular}[c]{@{}c@{}}$-20.1$\\{\scriptsize$[-24.5,-15.9]$}\end{tabular} & \begin{tabular}[c]{@{}c@{}}$-1.7$\\{\scriptsize$[-4.0,+0.5]$}\end{tabular} & 9.1\%; 25.2\% \\
 & 10 & \begin{tabular}[c]{@{}c@{}}$-48.8$\\{\scriptsize$[-53.2,-44.6]$}\end{tabular} & \begin{tabular}[c]{@{}c@{}}$-30.5$\\{\scriptsize$[-34.9,-26.2]$}\end{tabular} & --- & 9.1\%; 25.2\% \\
\bottomrule
\end{tabular}

\end{table}

\paragraph{A second reference draw.}
A pre-specified experiment redrew the reference for each of the 149 questions from the listed answers that neither contain nor are contained in the first (318 of the 596 answers change match group) and re-ran the referenced calls of all three graders, reusing their blind calls. All nine predictions were confirmed (\cref{tab:r22-redraw}). With the second reference, pooled accuracy is still 21.3, 16.9, and 17.1 points below the answer lists (Sonnet~5, Opus~5, Haiku~4.5), within 3 points of the first draw. Credit follows the reference for the answers whose match changes: correct answers that match the first reference but not the second go from 100\%, 99.3\%, and 99.3\% credited to 61.8\%, 72.1\%, and 74.3\%, and those that match only the second go from 64.4\%, 76.2\%, and 68.4\% to 99.4\%, 99.4\%, and 97.7\%. Reported without a prediction, each model's referenced accuracy moves by at most 8.1 points between draws; per-model tables are in the extended appendix.

\begin{table}[htbp]
\centering
\caption{The benchmark study under a second reference draw: pooled referenced minus answer-list accuracy under each draw and its change (points, 95\% CI, question bootstrap), and the share of correct answers credited (score 7 or above) among those that stop or start matching the reference between draws. Opus~5 uses the 136 questions with both draws complete.\pre}
\label{tab:r22-redraw}
\footnotesize
\setlength{\tabcolsep}{3pt}
\begin{tabular}{lrccccc}
\toprule
 & & \multicolumn{3}{c}{Referenced $-$ list (points)} & \multicolumn{2}{c}{Credited (\%), draw 1 $\to$ 2 ($n$)} \\
\cmidrule(lr){3-5}\cmidrule(lr){6-7}
Grader & $n$ & Draw 1 & Draw 2 & Change & Stop matching & Start matching \\
\midrule
Sonnet 5 & 149 & \begin{tabular}[c]{@{}c@{}}$-22.3$\\{\scriptsize$[-27.0,-17.9]$}\end{tabular} & \begin{tabular}[c]{@{}c@{}}$-21.3$\\{\scriptsize$[-25.8,-16.8]$}\end{tabular} & \begin{tabular}[c]{@{}c@{}}$+1.0$\\{\scriptsize$[-3.9,+5.9]$}\end{tabular} & 100 $\to$ 62 (144) & 64 $\to$ 99 (174) \\
Opus 5 & 136 & \begin{tabular}[c]{@{}c@{}}$-15.1$\\{\scriptsize$[-19.3,-11.0]$}\end{tabular} & \begin{tabular}[c]{@{}c@{}}$-16.9$\\{\scriptsize$[-21.3,-12.7]$}\end{tabular} & \begin{tabular}[c]{@{}c@{}}$-1.8$\\{\scriptsize$[-6.6,+2.9]$}\end{tabular} & 99 $\to$ 72 (140) & 76 $\to$ 99 (164) \\
Haiku 4.5 & 149 & \begin{tabular}[c]{@{}c@{}}$-20.1$\\{\scriptsize$[-24.5,-15.9]$}\end{tabular} & \begin{tabular}[c]{@{}c@{}}$-17.1$\\{\scriptsize$[-21.1,-13.3]$}\end{tabular} & \begin{tabular}[c]{@{}c@{}}$+3.0$\\{\scriptsize$[-1.3,+7.4]$}\end{tabular} & 99 $\to$ 74 (144) & 68 $\to$ 98 (174) \\
\bottomrule
\end{tabular}

\end{table}

\paragraph{A same-reference regrade.}
To tell a draw effect from call-to-call variation, a pre-specified experiment re-ran every referenced call of all three graders with the first reference (table in the extended appendix). All nine predictions were confirmed. Pooled referenced accuracy moves by less than 2 points between two calls with the same reference, and correct answers whose match changes between draws keep the credit of their first call to within about 3 points, so the credit they gain or lose under the second draw, 24 to 38 points beyond the regrade, is a draw effect. Gaps between models are not: the Haiku~4.5 minus Llama-3.1-8B gap under Haiku~4.5 as grader moves by 10.7 points under the same-reference regrade and by 10.7 under the second draw.
%==========================================================================================
\FloatBarrier
\section{Benchmark graders and fixes}
\label{app:r24}

\paragraph{Design.} A pre-specified experiment graded the 596 benchmark answers of \cref{sec:cost} (149 \textsc{MACE} questions, one answer each from Llama-3.1-8B-Instruct, Claude Haiku~4.5, Sonnet~5, and Opus~5; 90.6\% correct by the answer lists) with two graders copied verbatim from their repositories at pinned commits: the SimpleQA grader of \texttt{openai/simple-evals}, which sends a few-shot template to GPT-4.1 (system message ``You are a helpful assistant.'', temperature 0.5) and reads a letter A (correct), B, or C, and the judge prompt of Humanity's Last Exam, which BrowseComp also uses, sent to o3-mini with HLE's structured output or to GPT-4.1 as configured for BrowseComp and read as ``correct: yes/no''. Sonnet~5 (thinking disabled) ran both prompts as a second judge. Each answer was graded with the first reference draw as the gold target and with the full answer list written as ``Any one of the following is a complete correct answer: $a_1$; $a_2$; \dots''. Every call returned a parseable grade, and the official parse and a markdown-robust parse agreed on all of them (\cref{tab:r24-standard}).

\begin{table}[htbp]
\centering
\caption{Benchmark graders with one reference and with the full answer list: pooled accuracy minus answer-list accuracy (points, 95\% CI, question bootstrap), and the share of correct answers that differ from the reference (396) and of wrong answers (56) that each condition credits. 149 questions per row.\pre}
\label{tab:r24-standard}
\footnotesize
\setlength{\tabcolsep}{3pt}
\begin{adjustbox}{max width=\linewidth}\begin{tabular}{llcccccc}
\toprule
 & & \multicolumn{2}{c}{Accuracy minus list (points)} & \multicolumn{2}{c}{Differing correct credited (\%)} & \multicolumn{2}{c}{Wrong credited (\%)} \\
\cmidrule(lr){3-4}\cmidrule(lr){5-6}\cmidrule(lr){7-8}
Prompt & Judge & One reference & Full list & One ref. & Full list & One ref. & Full list \\
\midrule
SimpleQA grader & GPT-4.1 & \begin{tabular}[c]{@{}c@{}}$-30.4$\\{\scriptsize$[-35.4,-25.3]$}\end{tabular} & \begin{tabular}[c]{@{}c@{}}$+0.3$\\{\scriptsize$[-0.7,+1.5]$}\end{tabular} & 53.0 & 99.2 & 8.9 & 8.9 \\
BrowseComp judge & GPT-4.1 & \begin{tabular}[c]{@{}c@{}}$-66.8$\\{\scriptsize$[-71.8,-61.7]$}\end{tabular} & \begin{tabular}[c]{@{}c@{}}$-0.2$\\{\scriptsize$[-0.8,+0.5]$}\end{tabular} & 0.0 & 99.2 & 0.0 & 3.6 \\
HLE judge & o3-mini & \begin{tabular}[c]{@{}c@{}}$-66.6$\\{\scriptsize$[-71.8,-61.6]$}\end{tabular} & \begin{tabular}[c]{@{}c@{}}$+0.0$\\{\scriptsize$[-0.8,+0.8]$}\end{tabular} & 0.0 & 99.2 & 0.0 & 5.4 \\
SimpleQA grader & Sonnet~5 & \begin{tabular}[c]{@{}c@{}}$-34.4$\\{\scriptsize$[-39.9,-29.2]$}\end{tabular} & \begin{tabular}[c]{@{}c@{}}$+0.5$\\{\scriptsize$[-0.3,+1.5]$}\end{tabular} & 46.7 & 99.5 & 10.7 & 8.9 \\
HLE judge & Sonnet~5 & \begin{tabular}[c]{@{}c@{}}$-64.9$\\{\scriptsize$[-70.3,-59.9]$}\end{tabular} & \begin{tabular}[c]{@{}c@{}}$+0.2$\\{\scriptsize$[-0.3,+0.7]$}\end{tabular} & 2.3 & 99.7 & 1.8 & 3.6 \\
\bottomrule
\end{tabular}
\end{adjustbox}
\end{table}

The wrong answers credited only with the full list are ``The Democratic Republic of Congo'' (listed as ``Democratic Republic of the Congo'', credited in four configurations), ``Territory of New Guinea'' (listed ``Papua New Guinea''), ``Wales'' (listed ``United Kingdom''), and ``Sudan'' (listed ``South Sudan''); the registered bound of 3 points on extra wrong-answer credit failed for the BrowseComp/HLE judge with GPT-4.1 (2 of 56) and o3-mini (3 of 56).

\paragraph{GPT-4.1 with our rubric.} GPT-4.1, with the logprob-expected score of \cref{app:frontier}, graded the same answers blind, with the first reference, and with a second reference drawn from each list. Answers that stop matching under the second draw are credited 100\% under the first and 87.5\% under the second (144 answers); answers that start matching, 80.5\% and 99.4\% (174). All four predictions were confirmed (\cref{tab:r24-gpt41}).

\begin{table}[htbp]
\centering
\caption{GPT-4.1 as a benchmark grader with our rubric: accuracy (\%) by answering model, and pooled accuracy minus the answer lists (points, 95\% CI).\pre}
\label{tab:r24-gpt41}
\footnotesize
\begin{tabular}{lcccc}
\toprule
Model & Answer list & Blind & One reference & Second reference \\
\midrule
Llama-3.1-8B & 69.1 & 63.8 & 65.8 & --- \\
Haiku~4.5 & 94.6 & 94.6 & 80.5 & --- \\
Sonnet~5 & 98.7 & 94.0 & 79.9 & --- \\
Opus~5 & 100.0 & 100.0 & 98.0 & --- \\
Pooled, minus list & & \begin{tabular}[c]{@{}c@{}}$-2.5$\\{\scriptsize$[-4.7,-0.3]$}\end{tabular} & \begin{tabular}[c]{@{}c@{}}$-9.6$\\{\scriptsize$[-12.9,-6.4]$}\end{tabular} & \begin{tabular}[c]{@{}c@{}}$-7.0$\\{\scriptsize$[-10.1,-4.0]$}\end{tabular} \\
\bottomrule
\end{tabular}

\end{table}

\paragraph{The answer list and prompt changes, measured by grading.} On the 40 \textsc{MACE} subsamples of \cref{app:frontier}, Sonnet~5, Opus~5 (their earlier blind and referenced calls reused), and GPT-4.1 graded every candidate with the full answer list in the reference slot, and Sonnet~5 with adaptive thinking graded them with one reference and the full-credit sentence, its earlier thinking blind calls reused (\cref{tab:r24-fix}). All nine answer-list predictions and two of the three combination predictions were confirmed; the combination left no measurable rise in false rejections, so its third prediction, that the rise would remain, was not confirmed.

\begin{table}[htbp]
\centering
\caption{Grading AUROC of correct answers that differ from the reference against wrong answers, blind, with one reference, and with the full answer list, and changes against blind (95\% CI, task bootstrap); below, Sonnet~5 with the full-credit sentence, adaptive thinking, and both, each against its own blind calls. 40 tasks.\pre}
\label{tab:r24-fix}
\footnotesize
\setlength{\tabcolsep}{3pt}
\begin{adjustbox}{max width=\linewidth}\begin{tabular}{lcccccc}
\toprule
 & \multicolumn{3}{c}{Grading AUROC} & \multicolumn{3}{c}{Change against blind} \\
\cmidrule(lr){2-4}\cmidrule(lr){5-7}
Judge & Blind & One ref. & Full list & One ref. & Full list & False rejections, full list (pp) \\
\midrule
Sonnet~5 & 0.938 & 0.785 & 0.970 & \begin{tabular}[c]{@{}c@{}}$-0.153$\\{\scriptsize$[-0.219,-0.097]$}\end{tabular} & \begin{tabular}[c]{@{}c@{}}$+0.032$\\{\scriptsize$[+0.008,+0.060]$}\end{tabular} & \begin{tabular}[c]{@{}c@{}}$-2.2$\\{\scriptsize$[-8.0,+4.9]$}\end{tabular} \\
Opus~5 & 0.936 & 0.771 & 0.969 & \begin{tabular}[c]{@{}c@{}}$-0.164$\\{\scriptsize$[-0.223,-0.110]$}\end{tabular} & \begin{tabular}[c]{@{}c@{}}$+0.033$\\{\scriptsize$[-0.000,+0.084]$}\end{tabular} & \begin{tabular}[c]{@{}c@{}}$-6.3$\\{\scriptsize$[-12.8,-1.1]$}\end{tabular} \\
GPT-4.1 & 0.935 & 0.872 & 0.992 & \begin{tabular}[c]{@{}c@{}}$-0.063$\\{\scriptsize$[-0.133,+0.004]$}\end{tabular} & \begin{tabular}[c]{@{}c@{}}$+0.057$\\{\scriptsize$[+0.013,+0.121]$}\end{tabular} & \begin{tabular}[c]{@{}c@{}}$-2.2$\\{\scriptsize$[-4.6,-0.5]$}\end{tabular} \\
\midrule
\multicolumn{7}{@{}l}{\emph{Sonnet~5, when the answers are not listed: change against blind}}\\
 & \multicolumn{3}{l}{Setting} & AUROC & & False rejections (pp) \\
 & \multicolumn{3}{l}{Full-credit sentence} & \begin{tabular}[c]{@{}c@{}}$-0.065$\\{\scriptsize$[-0.114,-0.019]$}\end{tabular} & & \begin{tabular}[c]{@{}c@{}}$+18.1$\\{\scriptsize$[+9.9,+28.0]$}\end{tabular} \\
 & \multicolumn{3}{l}{Adaptive thinking} & \begin{tabular}[c]{@{}c@{}}$-0.091$\\{\scriptsize$[-0.140,-0.050]$}\end{tabular} & & \begin{tabular}[c]{@{}c@{}}$+14.9$\\{\scriptsize$[+8.2,+23.1]$}\end{tabular} \\
 & \multicolumn{3}{l}{Thinking and the sentence} & \begin{tabular}[c]{@{}c@{}}$-0.024$\\{\scriptsize$[-0.062,+0.006]$}\end{tabular} & & \begin{tabular}[c]{@{}c@{}}$-0.2$\\{\scriptsize$[-7.2,+4.6]$}\end{tabular} \\
\bottomrule
\end{tabular}
\end{adjustbox}
\end{table}

%==========================================================================================
\FloatBarrier
\section{Judges that reason before deciding}
\label{app:r25}

A pre-specified experiment graded the benchmark answers of \cref{sec:cost} with three prompts that ask the judge to reason about correctness, blind and with the first reference, with GPT-4.1 (temperature 0) and Sonnet~5 (thinking disabled), each allowed 1,024 output tokens; none of the prompts says whether a question can have several correct answers. \emph{Chain of thought}: ``Grade the response to the question, using the reference answer. [\ldots] Think step by step about whether the response is correct''. \emph{Verify first}: ``1.~Verify: using your own knowledge, check whether the response's final answer is correct for the question as asked. 2.~Compare: check the response against the reference answer. 3.~Decide whether the response is correct.'' \emph{Solve first}: ``First answer the question yourself. Then grade the response, using your own answer and the reference answer.'' Each ends with a line ``correct: yes'' or ``correct: no''; every one of the 7,152 replies had one, and none was cut off (\cref{tab:r25}). We had predicted a drop of at least 10 points in credit for differing correct answers; the drop was 0.8 to 8.6, so that prediction failed for all six.

\begin{table}[htbp]
\centering
\caption{Judges that reason before deciding, on 149 benchmark questions: pooled accuracy minus answer-list accuracy (points, 95\% CI), and the share of differing correct (396), matching correct (144), and wrong (56) answers credited blind and with the reference.\pre}
\label{tab:r25}
\footnotesize
\setlength{\tabcolsep}{3pt}
\begin{adjustbox}{max width=\linewidth}\begin{tabular}{llccccc}
\toprule
 & & \multicolumn{2}{c}{Accuracy minus list (points)} & \multicolumn{3}{c}{Credited (\%): blind $\to$ reference} \\
\cmidrule(lr){3-4}\cmidrule(lr){5-7}
Judge type & Judge & Blind & Reference & Differing correct & Matching correct & Wrong \\
\midrule
Chain of thought & GPT-4.1 & \begin{tabular}[c]{@{}c@{}}$-0.8$\\{\scriptsize$[-2.3,+0.7]$}\end{tabular} & \begin{tabular}[c]{@{}c@{}}$-3.0$\\{\scriptsize$[-4.7,-1.3]$}\end{tabular} & 97.7 $\to$ 94.7 & 98.6 $\to$ 98.6 & 10.7 $\to$ 8.9 \\
 & Sonnet~5 & \begin{tabular}[c]{@{}c@{}}$-0.5$\\{\scriptsize$[-1.8,+1.0]$}\end{tabular} & \begin{tabular}[c]{@{}c@{}}$-6.0$\\{\scriptsize$[-8.7,-3.5]$}\end{tabular} & 98.0 $\to$ 90.2 & 99.3 $\to$ 99.3 & 10.7 $\to$ 7.1 \\
Verify first & GPT-4.1 & \begin{tabular}[c]{@{}c@{}}$-0.3$\\{\scriptsize$[-1.8,+1.2]$}\end{tabular} & \begin{tabular}[c]{@{}c@{}}$-1.5$\\{\scriptsize$[-3.0,+0.0]$}\end{tabular} & 98.0 $\to$ 96.7 & 99.3 $\to$ 98.6 & 12.5 $\to$ 10.7 \\
 & Sonnet~5 & \begin{tabular}[c]{@{}c@{}}$-0.2$\\{\scriptsize$[-1.3,+1.0]$}\end{tabular} & \begin{tabular}[c]{@{}c@{}}$-6.2$\\{\scriptsize$[-8.9,-3.7]$}\end{tabular} & 98.7 $\to$ 90.2 & 99.3 $\to$ 99.3 & 8.9 $\to$ 5.4 \\
Solve first & GPT-4.1 & \begin{tabular}[c]{@{}c@{}}$-0.3$\\{\scriptsize$[-1.7,+1.0]$}\end{tabular} & \begin{tabular}[c]{@{}c@{}}$-0.5$\\{\scriptsize$[-2.0,+1.2]$}\end{tabular} & 98.2 $\to$ 97.5 & 98.6 $\to$ 100.0 & 12.5 $\to$ 12.5 \\
 & Sonnet~5 & \begin{tabular}[c]{@{}c@{}}$+0.0$\\{\scriptsize$[-1.3,+1.3]$}\end{tabular} & \begin{tabular}[c]{@{}c@{}}$-3.4$\\{\scriptsize$[-5.7,-1.2]$}\end{tabular} & 98.2 $\to$ 93.4 & 100.0 $\to$ 100.0 & 12.5 $\to$ 10.7 \\
\bottomrule
\end{tabular}
\end{adjustbox}
\end{table}

%==========================================================================================
\FloatBarrier
\section{Instruction against reasoning, grading with wrong answers, and RL verifiers}
\label{app:r26}

\paragraph{Instruction against reasoning.} A pre-specified experiment separated the instruction from written reasoning on the benchmark answers of \cref{sec:cost}. GPT-4.1 (temperature 0) and Sonnet~5 (thinking disabled) graded every answer with a direct prompt, ``Is the response a correct answer to the question?'', asking for a single ``correct: yes'' or ``correct: no'' line, blind and with the reference; with the HLE judge prompt verbatim; and with the same prompt after deleting ``do not argue for any answer different than [correct\_answer], focus only on whether the answers match'' and adding ``A response that gives a different answer that is also correct for the question is correct'' (\cref{tab:r26-instruction}). All settings were the same across prompts, so the verbatim rows differ slightly from \cref{tab:r24-standard}, where GPT-4.1 ran as simple-evals configures it. Without written reasoning, accuracy with the reference falls 17.4 points below the lists for GPT-4.1 and 10.1 for Sonnet~5 (12.1 and 8.4 below blind), against 3.0 and 6.0 with chain of thought; the registered prediction that the two would be within 5 points held for Sonnet~5 and failed for GPT-4.1. Removing the match clause recovers 35.4 and 52.5 points but leaves 31.2 and 12.6, so the registered 10-point bound failed for both.

\begin{table}[htbp]
\centering
\caption{Instruction against reasoning on the benchmark pool: accuracy minus answer-list accuracy (points, 95\% CI), and the share of differing correct (396) and wrong (56) answers credited. 149 questions.\pre}
\label{tab:r26-instruction}
\footnotesize
\setlength{\tabcolsep}{3pt}
\begin{adjustbox}{max width=\linewidth}\begin{tabular}{llcccc}
\toprule
Judge & Prompt & Condition & Accuracy minus list (points) & Differing correct credited (\%) & Wrong credited (\%) \\
\midrule
GPT-4.1 & Direct yes/no & blind & \begin{tabular}[c]{@{}c@{}}$-5.4$\\{\scriptsize$[-7.7,-3.2]$}\end{tabular} & 91.4 & 12.5 \\
 & Direct yes/no & reference & \begin{tabular}[c]{@{}c@{}}$-17.4$\\{\scriptsize$[-21.0,-14.1]$}\end{tabular} & 73.2 & 7.1 \\
 & HLE judge, verbatim & reference & \begin{tabular}[c]{@{}c@{}}$-66.6$\\{\scriptsize$[-71.8,-61.6]$}\end{tabular} & 0.0 & 0.0 \\
 & HLE judge, match clause removed & reference & \begin{tabular}[c]{@{}c@{}}$-31.2$\\{\scriptsize$[-36.4,-26.3]$}\end{tabular} & 53.3 & 0.0 \\
\addlinespace
Sonnet~5 & Direct yes/no & blind & \begin{tabular}[c]{@{}c@{}}$-1.7$\\{\scriptsize$[-3.5,+0.0]$}\end{tabular} & 96.5 & 10.7 \\
 & Direct yes/no & reference & \begin{tabular}[c]{@{}c@{}}$-10.1$\\{\scriptsize$[-13.3,-6.9]$}\end{tabular} & 83.3 & 10.7 \\
 & HLE judge, verbatim & reference & \begin{tabular}[c]{@{}c@{}}$-65.1$\\{\scriptsize$[-70.3,-60.1]$}\end{tabular} & 1.8 & 3.6 \\
 & HLE judge, match clause removed & reference & \begin{tabular}[c]{@{}c@{}}$-12.6$\\{\scriptsize$[-16.3,-9.1]$}\end{tabular} & 80.3 & 7.1 \\
\bottomrule
\end{tabular}
\end{adjustbox}
\end{table}

\paragraph{Grading with many wrong answers.} On the 40 \textsc{MACE} subsamples, the same judges graded every candidate with the direct, verify-first, and solve-first prompts, blind and with the reference, and with the verbatim and edited HLE judge (\cref{tab:r26-subsample}). With the reference, grading AUROC falls 0.04 to 0.09 under every prompt; only solve-first stays within the registered 0.05 for both judges. False acceptances rise by at most 3.2 points. The verbatim HLE judge rejects 96\% to 100\% of differing correct answers, at chance AUROC; removing its match clause brings AUROC to 0.61 and 0.82. For comparison, the full answer list puts AUROC 0.03 to 0.06 above blind (\cref{tab:r24-fix}).

\begin{table}[htbp]
\centering
\caption{Grading on the 40-task \textsc{MACE} subsample: grading AUROC of differing correct answers against wrong answers, blind and with the reference, and changes in AUROC, false rejections, and false acceptances (95\% CI, task bootstrap); last row per judge: the HLE judge with its match clause removed against the verbatim judge.\pre}
\label{tab:r26-subsample}
\footnotesize
\setlength{\tabcolsep}{3pt}
\begin{adjustbox}{max width=\linewidth}\begin{tabular}{llcccc}
\toprule
Judge & Prompt & AUROC, blind $\to$ reference & AUROC change & False rejections change (pp) & False acceptances change (pp) \\
\midrule
GPT-4.1 & Direct yes/no & 0.848 $\to$ 0.759 & \begin{tabular}[c]{@{}c@{}}$-0.089$\\{\scriptsize$[-0.141,-0.038]$}\end{tabular} & \begin{tabular}[c]{@{}c@{}}$+23.0$\\{\scriptsize$[+12.3,+34.3]$}\end{tabular} & \begin{tabular}[c]{@{}c@{}}$-5.1$\\{\scriptsize$[-11.2,+0.6]$}\end{tabular} \\
 & Verify first & 0.927 $\to$ 0.859 & \begin{tabular}[c]{@{}c@{}}$-0.068$\\{\scriptsize$[-0.115,-0.026]$}\end{tabular} & \begin{tabular}[c]{@{}c@{}}$+10.3$\\{\scriptsize$[+2.7,+19.5]$}\end{tabular} & \begin{tabular}[c]{@{}c@{}}$+3.2$\\{\scriptsize$[-0.2,+7.9]$}\end{tabular} \\
 & Solve first & 0.907 $\to$ 0.868 & \begin{tabular}[c]{@{}c@{}}$-0.040$\\{\scriptsize$[-0.088,-0.002]$}\end{tabular} & \begin{tabular}[c]{@{}c@{}}$+4.7$\\{\scriptsize$[+0.4,+10.8]$}\end{tabular} & \begin{tabular}[c]{@{}c@{}}$+3.2$\\{\scriptsize$[-2.0,+9.2]$}\end{tabular} \\
 & HLE edited minus verbatim & 0.485 $\to$ 0.611 & \begin{tabular}[c]{@{}c@{}}$+0.126$\\{\scriptsize$[+0.077,+0.176]$}\end{tabular} & \begin{tabular}[c]{@{}c@{}}$-28.5$\\{\scriptsize$[-39.1,-18.1]$}\end{tabular} & \begin{tabular}[c]{@{}c@{}}$+3.4$\\{\scriptsize$[+0.0,+8.7]$}\end{tabular} \\
\addlinespace
Sonnet~5 & Direct yes/no & 0.862 $\to$ 0.805 & \begin{tabular}[c]{@{}c@{}}$-0.057$\\{\scriptsize$[-0.106,-0.014]$}\end{tabular} & \begin{tabular}[c]{@{}c@{}}$+13.8$\\{\scriptsize$[+5.8,+23.6]$}\end{tabular} & \begin{tabular}[c]{@{}c@{}}$-2.5$\\{\scriptsize$[-8.9,+3.6]$}\end{tabular} \\
 & Verify first & 0.908 $\to$ 0.818 & \begin{tabular}[c]{@{}c@{}}$-0.090$\\{\scriptsize$[-0.143,-0.041]$}\end{tabular} & \begin{tabular}[c]{@{}c@{}}$+17.1$\\{\scriptsize$[+8.3,+26.8]$}\end{tabular} & \begin{tabular}[c]{@{}c@{}}$+0.9$\\{\scriptsize$[-3.2,+4.9]$}\end{tabular} \\
 & Solve first & 0.913 $\to$ 0.867 & \begin{tabular}[c]{@{}c@{}}$-0.047$\\{\scriptsize$[-0.082,-0.016]$}\end{tabular} & \begin{tabular}[c]{@{}c@{}}$+7.4$\\{\scriptsize$[+1.7,+14.7]$}\end{tabular} & \begin{tabular}[c]{@{}c@{}}$+1.9$\\{\scriptsize$[+0.0,+5.0]$}\end{tabular} \\
 & HLE edited minus verbatim & 0.501 $\to$ 0.818 & \begin{tabular}[c]{@{}c@{}}$+0.317$\\{\scriptsize$[+0.263,+0.369]$}\end{tabular} & \begin{tabular}[c]{@{}c@{}}$-65.9$\\{\scriptsize$[-76.5,-55.2]$}\end{tabular} & \begin{tabular}[c]{@{}c@{}}$+2.5$\\{\scriptsize$[+0.3,+5.7]$}\end{tabular} \\
\bottomrule
\end{tabular}
\end{adjustbox}
\end{table}

\paragraph{Open RL-reward verifiers.} A pre-specified experiment ran two open verifiers used for RL rewards on the benchmark answers, with one reference and with the full answer list as the gold answer: General-Reasoner's \texttt{general-verifier} (1.5B) with its template and reward rule (``Final Decision: Yes''), and CompassVerifier-3B with its template (\cref{tab:r26-verifiers}). Both ran locally with greedy decoding; after three low-memory stops the runs continued in chunks (Amendment 3); to answer a status question we counted credited answers in the partial General-Reasoner output (1,011 of 1,192 generations) before the run finished, without changing any setting (registration note); and 12 General-Reasoner replies reached 512 tokens without a decision and count as incorrect. With one reference, General-Reasoner's verifier falls 62.4 points short of the lists, crediting 5.8\% of differing correct answers and 97.2\% of matching ones. With the full list it credits 67.4\% of differing correct answers but only 77.8\% of matching ones, and accepts 30.4\% of wrong answers against 8.9\%, so it stays 24.2 points short and the registered 5-point bound failed. CompassVerifier-3B, whose prompt declares the standard answer ``ALWAYS CORRECT'', credits no differing correct answer and no wrong answer with one reference (66.6 points short), and with the full list credits 99.2\% of differing correct answers and 3.6\% of wrong ones, within 0.2 points of the lists.

\begin{table}[htbp]
\centering
\caption{Open RL-reward verifiers on the benchmark answers: accuracy minus answer-list accuracy (points, 95\% CI), and the share of differing correct (396) and wrong (56) answers credited, with one reference and with the full answer list. 149 questions.\pre}
\label{tab:r26-verifiers}
\footnotesize
\setlength{\tabcolsep}{3pt}
\begin{adjustbox}{max width=\linewidth}\begin{tabular}{lcccccc}
\toprule
 & \multicolumn{2}{c}{Accuracy minus list (points)} & \multicolumn{2}{c}{Differing correct credited (\%)} & \multicolumn{2}{c}{Wrong credited (\%)} \\
\cmidrule(lr){2-3}\cmidrule(lr){4-5}\cmidrule(lr){6-7}
Verifier & One reference & Full list & One ref. & Full list & One ref. & Full list \\
\midrule
General-Reasoner verifier (1.5B) & \begin{tabular}[c]{@{}c@{}}$-62.4$\\{\scriptsize$[-67.4,-57.4]$}\end{tabular} & \begin{tabular}[c]{@{}c@{}}$-24.2$\\{\scriptsize$[-28.7,-20.0]$}\end{tabular} & 5.8 & 67.4 & 8.9 & 30.4 \\
CompassVerifier-3B & \begin{tabular}[c]{@{}c@{}}$-66.6$\\{\scriptsize$[-71.8,-61.6]$}\end{tabular} & \begin{tabular}[c]{@{}c@{}}$-0.2$\\{\scriptsize$[-0.8,+0.5]$}\end{tabular} & 0.0 & 99.2 & 0.0 & 3.6 \\
\bottomrule
\end{tabular}
\end{adjustbox}
\end{table}

%==========================================================================================
\FloatBarrier
\section{P3 under correctness prompts, and hard single-answer questions}
\label{app:r27}

\paragraph{P3 under correctness prompts.} A pre-specified re-analysis of the stored round-25 and round-26 calls on the benchmark answers asked, for each prompt that asks whether the answer is correct, how many correct answers credited blind lose credit once the reference is shown (\cref{tab:r27-p3}). Answers that differ from the reference lose credit in every cell, from 1.3\% (solve first, GPT-4.1) to 22.4\% (direct prompt, GPT-4.1); answers that match it lose credit in at most 0.7\%. The difference excludes zero in seven of eight cells. These prompts have no same-prompt second blind call, so there is no noise floor.

\begin{table}[htbp]
\centering
\caption{Correct answers credited blind that lose credit with the reference, by prompt and judge (149 questions; differing and matching answers; difference with 95\% CI, question bootstrap).\pre}
\label{tab:r27-p3}
\footnotesize
\begin{tabular}{llccc}
\toprule
Judge & Prompt & Differing (\%) & Matching (\%) & Difference (points) \\
\midrule
GPT-4.1 & Direct yes/no & 22.4 (362) & 0.0 (139) & \begin{tabular}[c]{@{}c@{}}$+22.4$\\{\scriptsize$[+17.7,+27.0]$}\end{tabular} \\
 & Chain of thought & 4.4 (387) & 0.7 (142) & \begin{tabular}[c]{@{}c@{}}$+3.7$\\{\scriptsize$[+1.3,+6.2]$}\end{tabular} \\
 & Verify first & 2.1 (388) & 0.7 (143) & \begin{tabular}[c]{@{}c@{}}$+1.4$\\{\scriptsize$[-0.7,+3.2]$}\end{tabular} \\
 & Solve first & 1.3 (389) & 0.0 (142) & \begin{tabular}[c]{@{}c@{}}$+1.3$\\{\scriptsize$[+0.3,+2.5]$}\end{tabular} \\
Sonnet~5 & Direct yes/no & 14.1 (382) & 0.0 (142) & \begin{tabular}[c]{@{}c@{}}$+14.1$\\{\scriptsize$[+10.1,+18.4]$}\end{tabular} \\
 & Chain of thought & 8.5 (388) & 0.7 (143) & \begin{tabular}[c]{@{}c@{}}$+7.8$\\{\scriptsize$[+4.4,+11.4]$}\end{tabular} \\
 & Verify first & 9.2 (391) & 0.7 (143) & \begin{tabular}[c]{@{}c@{}}$+8.5$\\{\scriptsize$[+4.7,+12.7]$}\end{tabular} \\
 & Solve first & 4.9 (389) & 0.0 (144) & \begin{tabular}[c]{@{}c@{}}$+4.9$\\{\scriptsize$[+2.6,+7.5]$}\end{tabular} \\
\bottomrule
\end{tabular}

\end{table}

\paragraph{Hard single-answer questions.} The prompts that ask whether an answer is correct were tested above on questions the judges mostly answer correctly. A pre-specified experiment asked what happens on questions beyond the judge, where the reference is its main evidence. GPT-4.1-mini and Haiku~4.5 answered the 200 text-only exact-answer HLE questions of the second prevalence audit (\cref{app:prevalence}) with HLE's official system prompt; HLE's official judge (its verbatim prompt with o3-mini and structured output) labeled the 400 answers against the HLE answer, and only 7 were labeled correct. GPT-4.1 and Sonnet~5 then graded every answer with the verbatim HLE judge (reference only) and with the direct and solve-first prompts of \cref{app:r26}, blind and with the reference. Scored questions are those whose two answers have a label and a grade in every condition (191 for GPT-4.1, 196 for Sonnet~5).

With the reference, GPT-4.1 credited 17.3\% of the wrong answers under the direct prompt and 25.9\% when it solved first, against 0\% for the verbatim HLE judge; Sonnet~5 credited 0.5\% and 1.0\% (\cref{tab:r27b}). Two of the registered bounds on this increase failed for GPT-4.1, and a recall bound failed for Sonnet~5 by one of seven answers. Blind, both prompts credited far more wrong answers (GPT-4.1 61\% to 70\%; Sonnet~5 9\% to 21\%), so on these questions the reference is the judges' main evidence, and without the instruction to match GPT-4.1, though not Sonnet~5, often overrules it. Recall rests on 7 correct answers and is not informative.

Three caveats. (i) At the registered limit of 1,024 output tokens, solve-first replies often ended before a verdict (161 of GPT-4.1's 800 calls; 342 of the 460 Sonnet~5 solve-first calls run at that limit before the amendment); a dated amendment reran them at 4,096 tokens before any outcome was computed, after which 11 and 5 calls stayed ungraded. (ii) The labels come from a judge told to match, so some credited ``wrong'' answers may be equivalent to the reference. Post hoc, restricting to wrong answers that state an answer line leaves GPT-4.1 at 15.2\% (direct) and 22.9\% (solve first); at least 40 of its 97 solve-first false acceptances end with a reply saying the reference is wrong, and in a random 10 of GPT-4.1's solve-first false acceptances that we read, most credit a different final answer after the judge's own solution disagreed with the reference. GPT-4.1 credited more wrong answers from GPT-4.1-mini than from Haiku~4.5 (direct 20.7\% against 13.9\%; solve first 35.1\% against 16.6\%), so preference for its own family may contribute, but it overruled the reference for both.\ph{} (iii) Two graders and one benchmark.

\begin{table}[htbp]
\centering
\caption{Hard single-answer questions: share of label-wrong answers credited (95\% CI, question bootstrap) and label-correct answers credited, by grader and condition.\pre}
\label{tab:r27b}
\footnotesize
\begin{tabular}{lcccc}
\toprule
 & \multicolumn{2}{c}{GPT-4.1} & \multicolumn{2}{c}{Sonnet~5} \\
\cmidrule(lr){2-3}\cmidrule(lr){4-5}
Grader & Wrong credited (\%) & Recall & Wrong credited (\%) & Recall \\
\midrule
HLE judge, verbatim (reference) & 0.0 \ci{0.0}{0.0} & 6/7 & 0.0 \ci{0.0}{0.0} & 7/7 \\
Direct, blind & 61.1 \ci{55.9}{66.1} & 3/7 & 8.8 \ci{5.5}{12.4} & 1/7 \\
Direct, reference & 17.3 \ci{13.4}{21.4} & 7/7 & 0.5 \ci{0.0}{1.3} & 7/7 \\
Solve first, blind & 69.6 \ci{64.1}{74.7} & 6/7 & 21.0 \ci{16.6}{25.9} & 3/7 \\
Solve first, reference & 25.9 \ci{21.2}{30.6} & 7/7 & 1.0 \ci{0.3}{2.1} & 6/7 \\
\bottomrule
\end{tabular}

\end{table}

%==========================================================================================
\FloatBarrier
\section{The swap under a prompt that does not instruct comparison, and a chosen reference}
\label{app:r30}

Two registered experiments answer objections the earlier text conceded without testing: that our effects
might be prompt-following, since every headline number comes from a prompt that tells the judge to use or
compare with the reference, and that the size of the benchmark shortfall is a property of drawing the
reference at random.

\paragraph{The swap under the direct prompt.}
The direct prompt of \cref{app:r26} names the reference and says nothing about what to do with it; it is the
only prompt we have with that property. We re-ran the second reference draw of \cref{app:acceptance} under
it, so the candidate text is byte-identical and only the reference moves. The draw-1 arm is the existing
run of \cref{app:r26}; only the draw-2 calls are new. Among the 318 of 596 answers whose match group changes
between draws, credit follows the reference (149 questions, 596 calls per judge, no errors):

\begin{center}\footnotesize
\begin{tabular}{llccc}
\toprule
Judge & Class ($n$) & Draw 1 & Draw 2 & Change \\
\midrule
GPT-4.1 & stop matching (144) & 0.986 & 0.722 & $-0.264$ \ci{-0.352}{-0.178} \\
 & start matching (174) & 0.718 & 0.966 & $+0.247$ \ci{+0.164}{+0.328} \\
 & matches neither (222) & 0.743 & 0.712 & $-0.032$ \ci{-0.089}{+0.027} \\
\addlinespace[2pt]
Sonnet~5 & stop matching (144) & 1.000 & 0.868 & $-0.132$ \ci{-0.198}{-0.076} \\
 & start matching (174) & 0.810 & 0.989 & $+0.178$ \ci{+0.116}{+0.248} \\
 & matches neither (222) & 0.851 & 0.793 & $-0.059$ \ci{-0.107}{-0.011} \\
\bottomrule
\end{tabular}
\end{center}

\noindent Both judges confirm that answers losing the match lose credit and answers gaining it gain credit.
Two registered predictions failed. We predicted the shift would be smaller than under our rubric, which it
is for Sonnet~5 ($-0.132$ and $+0.178$ against $-0.382$ and $+0.350$) but not for GPT-4.1, whose shift is
larger than its rubric shift of $-0.125$ and $+0.189$ (\textbf{falsified}); the two are not
instrument-matched, since the rubric is a 0--10 score thresholded at 7 and the direct prompt is a binary
yes or no, so the comparison bounds the instruction's contribution only loosely. We also predicted that
answers matching neither reference would move by at most 5 points, which holds for GPT-4.1 but not for
Sonnet~5, whose placebo moves $-5.9$ points with an interval excluding zero (\textbf{falsified}); part of
Sonnet~5's movement is therefore not specific to the reference.

\paragraph{A chosen reference instead of a drawn one.}
We rebuilt the benchmark pool with the reference chosen by two pre-specified rules: \emph{first-listed},
\textsc{MACE}'s own ordering, which uses nothing about our models, and \emph{modal}, the listed answer most
of the four answering models produce, which is the most favourable reference available. Answers, labels and
lists are unchanged; only the reference and the match groups differ.

\begin{center}\footnotesize
\setlength{\tabcolsep}{4pt}
\begin{tabular}{llccc}
\toprule
Grader & Rule & Chosen $-$ list & Chosen $-$ random & Differing correct \\
\midrule
SimpleQA, GPT-4.1 & first-listed & $-0.258$ & $+0.045$ \ci{-0.003}{+0.092} & 0.519 \\
 & modal & $-0.091$ & $+0.213$ \ci{+0.164}{+0.263} & 0.504 \\
BrowseComp/HLE, o3-mini & first-listed & $-0.562$ & $+0.104$ \ci{+0.035}{+0.173} & 0.000 \\
 & modal & $-0.203$ & $+0.463$ \ci{+0.406}{+0.522} & 0.000 \\
Our rubric, Sonnet~5$^{\ast}$ & first-listed & $-0.181$ & $+0.042$ \ci{+0.005}{+0.079} & 0.663 \\
 & modal & $-0.079$ & $+0.144$ \ci{+0.104}{+0.188} & 0.595 \\
\bottomrule
\end{tabular}
\end{center}

\noindent{\footnotesize The last column is the share of correct answers differing from the reference that the
grader credits. $^{\ast}$Sonnet~5 grades its own answers here, among those of the four answering models.
Under the random draw the shortfall is $-0.304$ for the SimpleQA grader, $-0.666$ for the BrowseComp/HLE
judge and $-0.223$ for our rubric, and the credited share is 0.530, 0.000 and 0.652 respectively.}

\noindent A chosen reference shrinks the shortfall, by 21.3 points for the SimpleQA grader and 46.3 for the
HLE judge under the modal rule, and every interval on that difference excludes zero except the SimpleQA
grader's under the first-listed rule, where it includes zero (\textbf{split}). It never removes the
shortfall: every grader stays below its answer-list accuracy under both rules.

The narrowing is largely mechanical rather than better grading, and the HLE judge shows this in its purest
form. It credits 0.000 of the correct answers that differ from its reference under every rule, so its gap
falls only because fewer answers count as differing, 121 under the modal rule against 396 under the random
draw, not because it credits any of them. Choosing a canonical reference does not repair a grader that is
told to match; it makes fewer answers differ. Across all six cells the credited share moves by 0.0 to 5.7
points where the shortfall moves by 4.2 to 46.3. The two 0.0 cells are the HLE judge, which credits no
differing correct answer under either rule and so is pinned by the very effect at issue; among the graders
that credit anything the share moves by 1.1 to 5.7 points. This is why we report the
credited share as the statistic that moved least between our two reference rules, and the accuracy gap as
the one that depends on how the reference was picked.

\paragraph{Why a reference is given at all.}
\Cref{sec:fix} gives the reading that connects its two prescriptions; here we record why we state it as a
reading rather than a result. No experiment here manipulates judge capability while holding the number of
correct answers fixed, and the two settings we contrast, the multi-answer factual questions of
\cref{sec:cost} and the hard single-answer questions of \cref{app:r27}, differ in both at once.

%==========================================================================================
\FloatBarrier
\section{The swap under a prompt that forbids using the reference}
\label{app:r31}

Round 30 showed the swap survives a prompt that names the reference and says nothing about what to do with
it. It did not test the remaining form of the instruction-following objection: a prompt that tells the judge
\emph{not} to use the reference. This experiment does, and was registered before any call.

\paragraph{Design.} The template is byte-identical to the direct prompt of \cref{app:r26} except for one
inserted sentence: ``The reference answer is only one of several correct answers. Do not compare the response
to it, and do not treat agreeing or disagreeing with it as evidence; judge only whether the response is
itself a correct answer to the question.'' Any difference from the direct arm is therefore attributable to
the instruction rather than to rewording, though not to the instruction alone: both draws here are new
while the direct arm reuses its first draw, so call-to-call variation is not excluded. Both are new because
this template had never been run. Four
cells of 596 calls, no errors and no unparseable replies.

\begin{center}\footnotesize
\begin{tabular}{llccc}
\toprule
Judge & Class ($n$) & Draw 1 & Draw 2 & Change \\
\midrule
Sonnet~5 & stop matching (144) & 0.986 & 0.979 & $-0.007$ \ci{-0.031}{+0.015} \\
 & start matching (174) & 0.977 & 0.989 & $+0.012$ \ci{-0.011}{+0.034} \\
 & matches neither (222) & 0.968 & 0.946 & $-0.023$ \ci{-0.050}{+0.004} \\
\addlinespace[2pt]
GPT-4.1 & stop matching (144) & 0.993 & 0.972 & $-0.021$ \ci{-0.051}{+0.007} \\
 & start matching (174) & 0.977 & 0.977 & $+0.000$ \ci{-0.016}{+0.018} \\
 & matches neither (222) & 0.968 & 0.968 & $+0.000$ \ci{-0.012}{+0.013} \\
\bottomrule
\end{tabular}
\end{center}

\noindent The swap does not survive. Every interval includes zero, so both registered directional predictions
are \textbf{not confirmed} for both judges. Against the direct prompt on the same answers, GPT-4.1's
stop-matching shift falls from $-0.264$ to $-0.021$ and its start-matching shift from $+0.247$ to $+0.000$;
Sonnet~5's fall from $-0.132$ and $+0.178$ to $-0.007$ and $+0.012$. We registered that the shift would be
smaller than each judge's own shift under the direct prompt, and it is, by a wide margin (confirmed for both
judges in both directions); the placebo moves at most 2.3 points (confirmed).

\paragraph{The judges did not simply become permissive.} That is the reading this result has to survive, and
it does not hold: wrong answers are still rejected, while credit for correct answers rises sharply against
the direct prompt on the same pool.

\begin{center}\footnotesize
\begin{tabular}{lccc}
\toprule
Judge & Prompt & Correct credited & Wrong credited \\
\midrule
Sonnet~5 & direct & 0.878 & 0.107 \\
 & forbidden (draw 1 / 2) & 0.976 / 0.969 & 0.107 / 0.125 \\
\addlinespace[2pt]
GPT-4.1 & direct & 0.800 & 0.071 \\
 & forbidden (draw 1 / 2) & 0.978 / 0.972 & 0.196 / 0.214 \\
\bottomrule
\end{tabular}
\end{center}

\noindent Accuracy lands within 0.2 to 1.7 points of answer-list accuracy in every cell, against the 30 to 67
points by which the shipped graders fall short (\cref{sec:cost}). The cost differs by judge: close to free
for Sonnet~5, and a real rise in wrong answers credited for GPT-4.1.

\paragraph{What this does and does not settle.} It does not make the effect an artifact of our wording:
\cref{app:r30} already shows the swap is large under a prompt that names the reference and gives no
instruction about it. Read together, the two results bound the effect from both sides. It does not require an
instruction to compare, and for these two frontier judges it does not survive one forbidding it.
That instruction is not a general fix: on the four open judges it leaves a large effect on \textsc{MACE} ($-2.58$ for Qwen2.5-7B, $-1.78$ for Qwen2.5-14B) and removes none of it on the replication pool for Qwen2.5-7B (\cref{app:panel2}).

%==========================================================================================
\FloatBarrier
\section{Mitigations}
\label{app:mitigation}

\Cref{tab:protocols} summarizes the mitigations in grading and selection. The full tables behind it (score combinations, grading and selection by protocol, exact-scorer values, worst-case loss, judge coverage, reveal-call variants, and the score-resolution screen) are in the extended appendix.

\begin{table}[htbp]
\centering
\footnotesize
\caption{\textbf{What helps: a full-credit sentence reduces the grading harm; an exploratory max rule restores grading at a cost when the answer is unique.} Change against blind judging, except \emph{Reveal $-$ ref.}, which is the registered blind-then-reveal protocol (score blind, then rescore with the reference and an instruction to keep the blind score unless the reference shows the answer is wrong) minus the reference, both with the original scorer. \emph{Sentence}: the full-credit sentence of \cref{app:setup}. Bold: the Sentence and Max rule cells differ from the Reference cell, and Reveal $-$ ref.\ differs from zero (95\% intervals); ---: not run or not defined. Pre-specified: the API judges' Reference and Sentence values, Qwen2.5-7B's \textsc{MACE} Reference grading, and \textsc{MACE} Reference accuracy (the original scorer gives the same values); MATH-500 Reference accuracy is a re-analysis; every Reveal $-$ ref.\ and Max rule value, and the rest, is post hoc. Task counts are in the extended appendix.}
\label{tab:protocols}
\setlength{\tabcolsep}{3.1pt}
\begin{tabular}{@{}llrrrrrr@{}}
\toprule
& & \multicolumn{4}{c}{Grading AUROC} & \multicolumn{2}{c}{Accuracy (pp)} \\
\cmidrule(lr){3-6}\cmidrule(lr){7-8}
Pool & Judge & Reference & Sentence & Reveal $-$ ref. & Max rule & Reference & Max rule \\
\midrule
\multicolumn{8}{@{}l}{\emph{Several correct answers}}\\
\textsc{MACE} & Sonnet~5 & $-0.15$ & $\mathbf{-0.06}$ & --- & $\mathbf{-0.01}$ & --- & --- \\
 & Opus~5 & $-0.16$ & $\mathbf{-0.05}$ & --- & $\mathbf{-0.01}$ & --- & --- \\
 & Qwen2.5-7B & $-0.08$ & --- & $0.00$ & $\mathbf{-0.01}$ & $+2.0$ & $+4.0$ \\
 & phi-4 & $-0.06$ & --- & $+0.01$ & $\mathbf{-0.01}$ & $+6.1$ & $+6.1$ \\
$k$-pool, $k{>}1$ & Qwen2.5-7B & $-0.23$ & --- & $\mathbf{+0.11}$ & $\mathbf{+0.00}$ & $-30.3$ & $\mathbf{+3.0}$ \\
 & Qwen2.5-14B & $-0.17$ & --- & $\mathbf{+0.12}$ & $\mathbf{-0.03}$ & $-3.0$ & $+3.0$ \\
 & gemma-2-9b & $-0.19$ & --- & $\mathbf{-0.06}$ & $\mathbf{+0.00}$ & $-3.0$ & $+0.0$ \\
 & phi-4 & $-0.25$ & --- & $\mathbf{+0.04}$ & $\mathbf{-0.01}$ & $-15.2$ & $\mathbf{+3.0}$ \\
\multicolumn{8}{@{}l}{\emph{One correct answer}}\\
MATH-500, answer & Qwen2.5-7B & $+0.29$ & --- & --- & $\mathbf{+0.24}$ & $+28.6$ & $\mathbf{+23.8}$ \\
 & Qwen2.5-14B & $+0.25$ & --- & --- & $\mathbf{+0.22}$ & $+30.6$ & $+29.7$ \\
MATH-500, solution & Qwen2.5-7B & $+0.29$ & --- & $\mathbf{-0.04}$ & $\mathbf{+0.13}$ & $+30.5$ & $\mathbf{+23.8}$ \\
\bottomrule
\end{tabular}
\end{table}

\paragraph{Instructions and several references.}
On the synthetic grid at $k{=}16$, the any-valid-answer instruction closes phi-4's margin ($+1.91 \to +0.03$), shrinks Qwen2.5-14B's ($+3.34 \to +1.21$), and does little or nothing for smaller judges, even raising Qwen2.5-1.5B's on rubric~$a$ ($+3.58 \to +4.54$). Showing $\min(3,k)$ valid references removes the margin toward the original reference on the Qwen-generated grid without restoring blind-level inflation, and largely fails for Qwen2.5-7B on the Llama-generated grid.

\paragraph{Blind-then-reveal: a registered negative result.}
The protocol lowers inflation and the margin against the plain reference at $k{>}1$ (registered on OSF, confirmed), but it did not protect grading: unlike the max rule below, it does not keep grading AUROC within 0.03 of blind judging, except for Qwen2.5-14B ($-0.01$; extended appendix).\ph{} Its registered prediction that it keeps the $k{=}1$ benefit was not confirmed, and its registered accuracy predictions on \textsc{MACE}'s natural mix failed (\cref{app:reasoning}). On MATH-500, where the first call reproduces the frozen blind scores ($r=0.991$ and $0.973$), we predicted that it would keep less than half of the reference's benefit; it keeps nearly all of it, so that prediction was falsified. In inflation units, removing the displayed blind score from the second call restores the harm for all three judges tested. For Qwen2.5-7B, neutral wording costs $+0.51$ \ci{+0.28}{+0.73} of protection, and another judge's blind score gives $+0.38$ against $+0.01$ for its own.\pre

\paragraph{The max rule (post hoc).}
Keeping the higher of the blind and referenced scores keeps grading AUROC within 0.03 of blind judging, where showing the reference lowers it by 0.06 to 0.25, and keeps the reference's benefit on \textsc{MACE}; this protection is partly by construction, and on MATH-500 the rule trails the reference for Qwen2.5-7B.\ph{} No score combination is best in every pool. It is a diagnostic rather than a policy, and one property decides where it may be used: taking the higher score means the reference can never overturn a blind acceptance, so the rule keeps every wrong answer the blind judge already credited and its false acceptances are the union of those made under both conditions. That is safe only where the blind judge is itself selective. On the hard questions of \cref{app:r27} it is not --- blind GPT-4.1 credits 61\% to 70\% of wrong answers there --- so the rule would preserve all of them.

\paragraph{The validity gate.}
A second preregistered protocol multiplies the blind score by a separate call's probability that the candidate is valid, with the reference shown to that call as one valid example. For phi-4 it removes the inflation; for Qwen2.5-7B it raises the $k{=}2$ margin from $+5.02$ to $+6.62$, because the validity call itself anchors.


\section{A second judge panel: mitigations, and P1--P3 on two added judges}
\label{app:panel2}

This section reports a run completed after the main analyses (2026-09-18/19). Every number in it is
post hoc\ph{} unless named otherwise: no prediction was registered for it, and the totals of
\cref{tab:ledger-totals} are unchanged. Two judges are new to the panel, Qwen2.5-14B-Instruct and
gemma-2-9b-it; the others repeat earlier configurations and serve as reproduction checks.

\paragraph{Mitigations.}
\Cref{tab:panel2-mitig} grades the same candidates under four prompt-level fixes. The effect is the
paired change in a 0--10 score when a \emph{different} correct answer is shown as the reference,
on candidates an exact oracle marks correct, with an item-clustered bootstrap (10{,}000 resamples;
1{,}889 rows per \textsc{MACE} cell and 3{,}987 per replication cell). Three readings matter. The
fixes shrink the effect where it is large but none removes it: Qwen2.5-7B keeps $-1.22$ of 10 on
\textsc{MACE} under its best arm. They are pool-dependent: on the replication pool Qwen2.5-7B's
chain-of-thought and verify-first arms are indistinguishable from no fix at all (paired deltas
$+0.15$ \ci{-0.04}{+0.33} and $-0.04$ \ci{-0.23}{+0.14}). And verify-first, which asks the judge to
check the reference before scoring, is \emph{worse} than plain chain-of-thought on every judge
measured on \textsc{MACE} --- the opposite of the ordering we had expected, and the one prediction
of this run that a registration would have recorded as falsified. For GPT-4.1 and Sonnet~5 the
bare effect is only $-0.20$ and $-0.15$, so their arms are uninformative rather than negative.

\begin{table}[htbp]
\centering
\footnotesize
\caption{\textbf{Prompt-level fixes shrink the effect but do not remove it.} Paired change in a
0--10 score when a different correct answer is the reference (95\% intervals). \emph{None}: the bare
rubric. \emph{Several}: a sentence stating that several answers may be correct. \emph{Forbid}: an
instruction not to compare. \emph{CoT}: chain-of-thought. \emph{Verify}: verify-first. All post hoc.\ph}
\label{tab:panel2-mitig}
\setlength{\tabcolsep}{3pt}
\begin{adjustbox}{max width=\linewidth}
\begin{tabular}{@{}llllll@{}}
\toprule
Judge & None & Several & Forbid & CoT & Verify \\
\midrule
\multicolumn{6}{@{}l}{\emph{\textsc{MACE} (factual, several correct answers)}}\\
Qwen2.5-7B & $-3.42$ \ci{-3.83}{-3.03} & $-2.45$ \ci{-2.80}{-2.12} & $-2.58$ \ci{-2.90}{-2.26} & $-1.22$ \ci{-1.53}{-0.91} & $-2.06$ \ci{-2.41}{-1.71} \\
Qwen2.5-14B & $-3.00$ \ci{-3.29}{-2.70} & $-1.33$ \ci{-1.58}{-1.07} & $-1.78$ \ci{-2.07}{-1.48} & $-0.92$ \ci{-1.14}{-0.70} & $-0.93$ \ci{-1.22}{-0.65} \\
gemma-2-9b & $-2.13$ \ci{-2.52}{-1.77} & $-0.61$ \ci{-0.85}{-0.39} & $-0.11$ \ci{-0.21}{-0.02} & $-0.31$ \ci{-0.49}{-0.14} & $-0.73$ \ci{-1.01}{-0.47} \\
phi-4 & $-1.08$ \ci{-1.34}{-0.84} & $-0.09$ \ci{-0.24}{+0.06} & $-0.02$ \ci{-0.15}{+0.11} & $-0.07$ \ci{-0.20}{+0.07} & $-0.03$ \ci{-0.24}{+0.19} \\
GPT-4.1 & $-0.20$ \ci{-0.38}{-0.08} & $-0.01$ \ci{-0.10}{+0.07} & $-0.10$ \ci{-0.20}{-0.01} & $-0.01$ \ci{-0.06}{+0.03} & $-0.01$ \ci{-0.05}{+0.03} \\
Sonnet~5 & $-0.15$ \ci{-0.27}{-0.05} & $-0.13$ \ci{-0.23}{-0.04} & $-0.06$ \ci{-0.13}{+0.01} & $-0.12$ \ci{-0.23}{-0.03} & $-0.10$ \ci{-0.20}{-0.00} \\
\midrule
\multicolumn{6}{@{}l}{\emph{Replication pool (arithmetic, enumerable answers)}}\\
Qwen2.5-7B & $-1.15$ \ci{-1.26}{-1.04} & $-1.00$ \ci{-1.12}{-0.88} & $-1.05$ \ci{-1.18}{-0.91} & $-1.00$ \ci{-1.17}{-0.83} & $-1.19$ \ci{-1.37}{-1.02} \\
Qwen2.5-14B & $-1.22$ \ci{-1.34}{-1.10} & $-0.70$ \ci{-0.81}{-0.59} & $-1.07$ \ci{-1.18}{-0.97} & $-0.52$ \ci{-0.63}{-0.42} & $-0.41$ \ci{-0.49}{-0.34} \\
gemma-2-9b & $-1.31$ \ci{-1.46}{-1.17} & $-0.18$ \ci{-0.23}{-0.12} & $+0.10$ \ci{+0.03}{+0.17} & $-0.12$ \ci{-0.18}{-0.05} & $-0.12$ \ci{-0.19}{-0.06} \\
phi-4 & $-0.11$ \ci{-0.17}{-0.05} & $+0.04$ \ci{+0.01}{+0.08} & $-0.04$ \ci{-0.08}{-0.01} & $+0.06$ \ci{+0.01}{+0.11} & $-0.02$ \ci{-0.07}{+0.02} \\
\bottomrule
\end{tabular}
\end{adjustbox}
\end{table}

\paragraph{P1 on two added judges.}
We reran the frozen rotation analysis of \cref{app:rotation} with three judges added to its
registered pair. The registered cells reproduce exactly (Qwen2.5-7B $+4.97$, phi-4 $+4.09$ under
rubric~$a$), and the added judges behave the same way (\cref{tab:panel2-rot}): text the reference
moves onto gains, text it moves off loses about as much, and text matching neither reference stays
within $+0.31$ in all ten cells. The frozen script and its stored results were not modified; the
extension ran from a copy.

\begin{table}[htbp]
\centering
\footnotesize
\caption{\textbf{The swap, with three judges added to the registered pair.} Change in score of
identical text when the reference is rotated onto it, off it, or onto neither (0--10, 95\%
intervals). Qwen2.5-7B and phi-4 are the registered judges and reproduce their stored values;
Sonnet~5, Qwen2.5-14B, and gemma-2-9b are post hoc.\ph}
\label{tab:panel2-rot}
\setlength{\tabcolsep}{4pt}
\begin{tabular}{@{}llccc@{}}
\toprule
Judge & Rubric & Onto & Off & Neither \\
\midrule
Qwen2.5-7B\pre & $a$ & $+4.97$ \ci{+4.66}{+5.27} & $-5.04$ \ci{-5.38}{-4.67} & $+0.13$ \ci{-0.04}{+0.31} \\
 & $b$ & $+6.34$ \ci{+6.10}{+6.60} & $-6.58$ \ci{-6.89}{-6.26} & $+0.10$ \ci{-0.11}{+0.29} \\
phi-4\pre & $a$ & $+4.09$ \ci{+3.79}{+4.38} & $-3.92$ \ci{-4.26}{-3.59} & $+0.30$ \ci{+0.04}{+0.57} \\
 & $b$ & $+3.25$ \ci{+2.95}{+3.51} & $-3.18$ \ci{-3.46}{-2.87} & $+0.24$ \ci{+0.01}{+0.47} \\
Sonnet~5\ph & $a$ & $+3.04$ \ci{+2.48}{+3.60} & $-2.66$ \ci{-3.32}{-2.03} & $+0.04$ \ci{-0.39}{+0.48} \\
 & $b$ & $+2.54$ \ci{+2.07}{+3.00} & $-2.51$ \ci{-3.14}{-1.93} & $+0.29$ \ci{-0.11}{+0.70} \\
Qwen2.5-14B\ph & $a$ & $+3.65$ \ci{+3.37}{+3.92} & $-3.52$ \ci{-3.82}{-3.24} & $+0.27$ \ci{+0.09}{+0.46} \\
 & $b$ & $+4.82$ \ci{+4.37}{+5.24} & $-4.79$ \ci{-5.25}{-4.32} & $+0.21$ \ci{-0.10}{+0.51} \\
gemma-2-9b\ph & $a$ & $+5.54$ \ci{+5.24}{+5.83} & $-5.33$ \ci{-5.76}{-4.88} & $+0.28$ \ci{+0.01}{+0.56} \\
 & $b$ & $+5.24$ \ci{+4.97}{+5.51} & $-5.11$ \ci{-5.58}{-4.64} & $+0.21$ \ci{-0.03}{+0.45} \\
\bottomrule
\end{tabular}
\end{table}

\paragraph{P2 and P3 on two added judges.}
\Cref{tab:panel2-p2p3} gives both. For P2 we use the per-task $G_{nm}$ of \cref{app:grid}: the
AUROC separating correct answers that differ from the reference from wrong ones, referenced minus
blind. The two registered judges reproduce their stored replication values to four decimals
($-0.2924$ and $-0.2754$), which is the check that the reimplementation is faithful; the two added
judges fall about half as far, and every interval excludes zero. The \textsc{MACE} column has no
stored open-judge value to reproduce against, so its matches are recomputed with the rule of
\cref{app:mace} and it should be read as corroboration rather than replication. For P3 we count
correct answers the judge scored 10 with no reference, then ask how many fall to 3 or below once a
reference is shown. The asymmetry is the result: among answers that \emph{match} the reference,
not one of the 2{,}998 accepted answers was taken back, in any judge or pool.

\begin{table}[htbp]
\centering
\footnotesize
\caption{\textbf{P2 and P3 across the panel.} $G_{nm}$: grading AUROC for correct answers that
differ from the reference against wrong ones, referenced minus blind (95\% intervals). Reversal:
of correct answers scored 10 blind, the share that falls to 3 or below once a reference is shown,
for answers that differ from it and for answers that match it. Qwen2.5-7B and phi-4 reproduce
stored replication $G_{nm}$ values\pre; all other cells are post hoc.\ph}
\label{tab:panel2-p2p3}
\setlength{\tabcolsep}{4pt}
\begin{adjustbox}{max width=\linewidth}
\begin{tabular}{@{}lccrrrr@{}}
\toprule
& \multicolumn{2}{c}{$G_{nm}$ (ref.\ $-$ blind)} & \multicolumn{2}{c}{Reversal, \textsc{MACE}} & \multicolumn{2}{c}{Reversal, replication} \\
\cmidrule(lr){2-3}\cmidrule(lr){4-5}\cmidrule(lr){6-7}
Judge & Replication & \textsc{MACE} & Differs & Matches & Differs & Matches \\
\midrule
Qwen2.5-7B\pre{} & $-0.292$ \ci{-0.323}{-0.261} & $-0.110$ \ci{-0.153}{-0.065} & 80.5\% (33/41) & 0/44 & 0.0\% (0/421) & 0/421 \\
phi-4\pre{} & $-0.275$ \ci{-0.305}{-0.246} & $-0.092$ \ci{-0.123}{-0.060} & 5.2\% (13/252) & 0/252 & 0.1\% (1/764) & 0/764 \\
Qwen2.5-14B & $-0.158$ \ci{-0.187}{-0.131} & $-0.061$ \ci{-0.095}{-0.028} & 19.8\% (21/106) & 0/106 & 0.2\% (1/621) & 0/621 \\
gemma-2-9b & $-0.159$ \ci{-0.182}{-0.137} & $-0.042$ \ci{-0.079}{-0.006} & 15.8\% (44/278) & 0/278 & 4.1\% (29/707) & 0/707 \\
GPT-4.1 & --- & --- & 0.7\% (2/283) & 0/283 & 1.5\% (11/732) & 0/732 \\
Sonnet~5 & --- & --- & 1.0\% (3/286) & 0/286 & 1.0\% (8/822) & 0/822 \\
\bottomrule
\end{tabular}
\end{adjustbox}
\end{table}

\paragraph{What we do not claim.}
Across the four open judges the share of the effect a fix removes rises as the bare effect falls,
a rank correlation of $-1.00$ for both the several-answers sentence and chain-of-thought (exact
$p = 0.0417$, $n=4$). We do not report this as a finding. Adding GPT-4.1 and Sonnet~5 collapses it
to $-0.14$, because ``share removed'' is a ratio on a bare effect of $-0.20$ or $-0.15$ for them;
and the replication pool cannot test it at all, since its four judges span a 1.14$\times$ range of
bare severity against \textsc{MACE}'s 3.17$\times$. The gradient is a four-judge, one-pool
observation, and we leave it as one.


\paragraph{What this costs the benchmarks themselves.}
The 30 to 67 point shortfalls of \cref{sec:cost} are measured on \textsc{MACE} questions chosen
because several answers are correct. They are not shortfalls on the SimpleQA, BrowseComp, or HLE
leaderboards, where such questions are rare. To put a number on the latter we combine three measured
quantities: the share of a benchmark's own questions that admit several correct answers, from the
1{,}000-question audit of \cref{app:prevalence}; the share of correct model answers that differ from a
randomly drawn reference, 69.0\% (5{,}223 of 7{,}569) on \textsc{MACE} questions with $k>1$; and the
share of such answers each shipped grader denies credit, from \cref{app:r24}. \Cref{tab:panel2-impact}
gives the product, an upper bound on the accuracy a model loses, with the prevalence and
differing-answer uncertainties propagated by a Jeffreys
posterior on each and 20{,}000 draws.

The effect is small: at most 0.4 points for SimpleQA and 2 points for HLE, matching the figures quoted in the abstract.
These are ceilings rather than expectations, because the product conditions on the model already
answering the question correctly and does not multiply by how often it does; the true leaderboard
delta is smaller by that factor.
Two caveats cut in opposite directions. The differing-answer share is measured on \textsc{MACE} and
assumed to transfer, as is the denial rate, which \cref{app:r24} measures on the same 149 \textsc{MACE}
questions; neither can be checked on the native benchmarks without answer lists they do not publish. And the audit's \textsc{unclear} labels are not counted as multi-answer: were every one of
them to go the other way, HLE would reach 20 points and BrowseComp 37, since 108 of BrowseComp's 200
audited questions are \textsc{unclear}. BrowseComp's $0.00$ should therefore be read as a bound that its audit cannot tighten, not as an estimate. The honest
summary is that this is a large failure mode with a currently small leaderboard cost, and the cost
scales with how often the questions admit several answers --- which is why the RL training sets of
\cref{sec:cost}, at 5.5\% to 44\%, are the setting to watch.

\begin{table}[htbp]
\centering
\footnotesize
\caption{\textbf{Expected shortfall on the benchmarks' own leaderboards.} Accuracy points a model
loses to single-reference grading, as prevalence $\times$ differing-answer share $\times$ grader denial
rate, with 95\% intervals from the prevalence and differing-answer posteriors. Prevalence is from the
adjudicated audit of \cref{app:prevalence} (200 questions each). Contrast the 30 to 67 point gaps of
\cref{sec:cost}, which are measured on questions selected for having several answers. All post hoc.\ph}
\label{tab:panel2-impact}
\setlength{\tabcolsep}{5pt}
\begin{tabular}{@{}llccc@{}}
\toprule
Benchmark & Grader as shipped & Multi-answer & Denies credit & Expected shortfall \\
\midrule
SimpleQA & SimpleQA grader + GPT-4.1 & 1.0\% & 47.0\% & $0.32$ \ci{0.07}{1.02} \\
 & SimpleQA grader + Sonnet~5 & 1.0\% & 53.3\% & $0.37$ \ci{0.08}{1.16} \\
HLE & HLE judge + o3-mini & 3.0\% & 100.0\% & $2.07$ \ci{0.87}{4.21} \\
 & HLE judge + GPT-4.1 & 3.0\% & 100.0\% & $2.07$ \ci{0.87}{4.18} \\
 & HLE judge + Sonnet~5 & 3.0\% & 97.7\% & $2.02$ \ci{0.85}{4.11} \\
BrowseComp & BrowseComp/HLE judge + GPT-4.1 & 0.0\% & 100.0\% & $0.00$ \ci{0.00}{0.87} \\
 & BrowseComp/HLE judge + Sonnet~5 & 0.0\% & 97.7\% & $0.00$ \ci{0.00}{0.83} \\
\bottomrule
\end{tabular}
\end{table}


\paragraph{Five primary endpoints, corrected for multiplicity.}
The ledger records 329 registered predictions and applies no multiplicity correction, which makes
the confirmatory boundary hard to read. We therefore name the five endpoints that carry the paper's
three claims, all registered before their analyses, and correct across that family with the
Holm step-down procedure at $\alpha = 0.05$:
\begin{itemize}[leftmargin=1.5em,itemsep=1pt,topsep=2pt]
\item[\textbf{E1}] the swap gain, P-RR1, $+4.97$ \ci{+4.66}{+5.27} (Qwen2.5-7B, rubric~$a$); the
registration predicts the \emph{sign} in all four judge~$\times$~rubric cells, which hold at $+3.25$ to
$+6.34$, and we quote one cell's magnitude here;
\item[\textbf{E2}] the swap loss, P-RR2, $-5.04$ \ci{-5.38}{-4.67};
\item[\textbf{E3}] gain minus placebo, P-RR3, $+4.81$ \ci{+4.44}{+5.17}, the specificity contrast;
\item[\textbf{E4}] the fall in grading AUROC, P-R4D4, $-0.275$ \ci{-0.305}{-0.245}, as recorded
when that round was scored; deciding matches from the extracted answer instead of by the registered
rule gives $-0.292$ \ci{-0.323}{-0.261} on the same judge and pool (\cref{tab:labelings-gnm});
\item[\textbf{E5}] retraction against a second blind call: among 843 accepted answers, the referenced
call and the second blind call disagree on 115, and all 115 disagreements go against the referenced
call (exact two-sided binomial $p = 4.8\times10^{-35}$). The two answers the second blind call scored
$\le 3$ were also scored $\le 3$ with the reference, so they are concordant and do not enter the test.
\end{itemize}
All five survive Holm correction. E1 to E4 are interval endpoints: each is an item-clustered
bootstrap whose 95\% interval excludes zero, and no resample of the 1{,}000 crosses it, which bounds
each raw $p$ below $0.001$ and so clears the tightest threshold in the family, $\alpha/5 = 0.01$.
Only E5 admits an exact $p$, $4.8\times10^{-35}$, against its threshold of $0.05$. A 1{,}000-resample
bootstrap cannot resolve a $p$ smaller than about $10^{-3}$, so we report bounds for E1 to E4 rather
than point $p$-values (\texttt{papers/analysis/primary\_endpoints\_holm.json}). E1 to E3 come from the frozen rotation
replication and E4 from the registered labels, so none depends on the any-format labeling that
\cref{app:grid} shows moves effect sizes. The remaining ledger entries should be read as robustness
and exploration rather than as confirmatory tests, and the totals in \cref{tab:ledger-totals} are
reported in that spirit.


\section{A pre-registered holdout replication of the swap}
\label{app:holdout}

Every rotation result above comes from one 120-item pool. This section reports a replication on
items that did not exist when the manuscript was written, with hypotheses, thresholds and verdict
rule fixed in \texttt{registration\_round32.md} and the analysis frozen in
\texttt{papers/analysis/holdout\_analysis.py}, both committed before the first judging call. It was run
once and is reported as it came out.\pre{} The registration text is also archived at \texttt{osf.io/npd9u},
posted after the analysis and so carrying no pre-data claim of its own; the pre-data record is the commit.

\paragraph{Pool.} 110 constraint items from the same five families and the same builder, drawn on a
new seed and filtered so that no item shares a family and parameter dictionary with the original
pool: overlap is 0 of 110. \textsc{digitsum} contributes 14 items rather than 24 because a
2{,}000-item draw yields only 14 novel instances of that family; the imbalance was registered rather
than repaired by reusing seen items. Candidates are 64 samples per item from Llama-3.1-8B-Instruct,
oracle pass rate 39.5\%, with 76 of the 110 items inside the $[0.2, 0.8]$ band the original design
targeted. Because the weights repository is gated and the cluster holds no credential, the
candidates were sampled from an ungated mirror of the same release whose configuration matches on
every architectural field; that substitution is recorded in the registration.

\paragraph{Result.} All five registered hypotheses hold in all four judge~$\times$~rubric cells
(\cref{tab:holdout}). Text the reference moves onto gains 3.7 to 6.3 points, text it moves off loses
about as much, and the two contrasts that carry the causal reading behave as registered: the
gainer-minus-placebo contrast is positive everywhere, and the placebo itself stays within $0.10$ of
zero in every cell --- tighter than the original run's $+0.10$ to $+0.30$. Point estimates land
within 0.4 points of the corresponding frozen cells (\cref{tab:rotation}) despite sharing no items.

\begin{table}[htbp]
\centering
\footnotesize
\caption{\textbf{The swap replicates on a pre-registered holdout.} Change in score of identical text
when the reference is rotated onto it, off it, or onto neither (0--10, 95\% intervals, item-clustered
bootstrap). Registered as H32-1 to H32-5 before any judging call; all five confirmed in all four
cells. $n$ differs by column because a class enters an item only when that item contains a candidate
of it.\pre}
\label{tab:holdout}
\setlength{\tabcolsep}{4pt}
\begin{tabular}{@{}llcccc@{}}
\toprule
Judge & Rubric & Onto (H32-1, H32-5) & Off (H32-2) & Neither (H32-4) & Onto $-$ neither (H32-3) \\
\midrule
Qwen2.5-7B & $a$ & $+4.88$ \ci{+4.59}{+5.21} & $-5.17$ \ci{-5.54}{-4.80} & $-0.00$ \ci{-0.15}{+0.14} & $+4.87$ \ci{+4.52}{+5.20} \\
 & $b$ & $+6.32$ \ci{+6.02}{+6.62} & $-6.71$ \ci{-7.02}{-6.39} & $+0.09$ \ci{-0.03}{+0.23} & $+6.20$ \ci{+5.92}{+6.54} \\
phi-4 & $a$ & $+4.40$ \ci{+4.11}{+4.68} & $-4.11$ \ci{-4.44}{-3.78} & $+0.10$ \ci{-0.09}{+0.28} & $+4.27$ \ci{+3.96}{+4.58} \\
 & $b$ & $+3.67$ \ci{+3.39}{+3.97} & $-3.68$ \ci{-3.98}{-3.39} & $-0.00$ \ci{-0.15}{+0.18} & $+3.59$ \ci{+3.33}{+3.91} \\
\midrule
\multicolumn{6}{@{}l}{\emph{Items contributing: 109 onto, 93 off, 88 neither (of 110).}}\\
\bottomrule
\end{tabular}
\end{table}


\section{What a single-reference reward does to a policy, in simulation}
\label{app:rlsim}

\Cref{sec:cost} states the reward-signal risk as a concern rather than a demonstrated harm, because
our preference-training runs were inconclusive. This section adds a mechanism check that is honest
about its own limits.\ph{}

\paragraph{What this is.} A policy-gradient simulation over scores that already exist. Each
replication-pool task carries 64 stored candidate responses, an exact oracle label for each, whether
each matches the shown reference, and each candidate's expected judge score with and without the
reference. The policy is a softmax over those 64 fixed candidates, trained by \textsc{reinforce}
against one of three rewards: the exact oracle, the judge with the reference shown, and the judge
with no reference. The third is the control: any drift it produces comes from the judge's own
preferences rather than from the displayed answer. Seeds are fixed per item and reward, so the
numbers below reproduce exactly.

\paragraph{What this is not.} Evidence that RL training on a reference-guided verifier distorts a
real policy. A real policy explores an open output space; this one can only move probability mass
between answers that were already sampled. We report it as a mechanism check and leave the claim in
\cref{sec:cost} as it stands.

\paragraph{Result.} \Cref{tab:rlsim} gives all seven cells. Rewarded by the exact oracle, a
converged policy puts 15.3\% of its valid probability mass on the answer that matches the reference
on the arithmetic pool, and 37.9\% on \textsc{MACE} --- close to the rate at which that answer
appears among correct candidates. Rewarded by the same judge \emph{with the reference shown}, it
puts 61.6\% to 86.6\% there. Paired by item, that is a rise of $+0.32$ to $+0.54$ over the oracle
policy, and every interval excludes zero.

The blind-judge control separates the reference from the judge. On \textsc{MACE} it is null for all
five judges, between $-0.038$ and $-0.016$ with every interval covering zero: nothing but the
displayed answer moves the policy. On the arithmetic pool it is small but not zero ($+0.029$ and
$+0.041$), which we read as the reference there often being the obvious answer, which a blind judge
also prefers. Either way the referenced signal is an order of magnitude larger.

\paragraph{A judge can look clean on grading and still distort a policy.} Sonnet~5 is the case worth
attention. Its bare anchoring on \textsc{MACE} is $-0.15$ of 10, the smallest of any judge we score,
and none of the four prompt-level fixes moves it detectably because there is so little to move
(\cref{tab:panel2-mitig}). By the measurements an evaluation team would normally run, it is the
clean judge in this paper. As a reward signal it still carries 37.9\% to 70.2\% of the policy's
valid mass onto the shown answer, $+0.322$ \ci{+0.249}{+0.395}. Grading quality and reward-signal
safety are not the same property, and the second does not follow from the first.

Two further readings. Under either judge reward the share of mass on \emph{any} correct answer moves
together --- on \textsc{MACE}, 0.96 under the oracle against 0.65 to 0.94 under blind and referenced
alike --- so that loss is ordinary grading error, not the reference's doing. And the concentration is
what a reward-hacking account predicts: the policy is not becoming less correct, it is becoming
correct in one particular way --- the way the grader was shown.

\begin{table}[htbp]
\centering
\footnotesize
\caption{\textbf{A single-reference reward concentrates a policy on the answer it was shown.} Share
of converged valid probability mass sitting on the reference-matching answer, by reward signal, with
the paired per-item difference from the oracle policy (95\% intervals). \emph{Blind} is the control:
the same judge with no reference. Simulation over stored scores on a fixed 64-candidate set, not RL
training. All post hoc.\ph}
\label{tab:rlsim}
\setlength{\tabcolsep}{4pt}
\begin{tabular}{@{}llccccc@{}}
\toprule
Pool & Judge & Oracle & Blind & Referenced & Ref $-$ oracle & Blind $-$ oracle \\
\midrule
Arithmetic & Qwen2.5-7B & $0.153$ & $0.182$ & $0.694$ & $+0.541$ \ci{+0.470}{+0.615} & $+0.029$ \ci{+0.006}{+0.056} \\
 & phi-4 & $0.153$ & $0.194$ & $0.616$ & $+0.463$ \ci{+0.386}{+0.537} & $+0.041$ \ci{+0.021}{+0.067} \\
\midrule
\textsc{MACE} & Qwen2.5-7B & $0.379$ & $0.357$ & $0.828$ & $+0.448$ \ci{+0.375}{+0.523} & $-0.023$ \ci{-0.080}{+0.036} \\
 & Qwen2.5-14B & $0.379$ & $0.343$ & $0.777$ & $+0.398$ \ci{+0.322}{+0.468} & $-0.036$ \ci{-0.093}{+0.025} \\
 & gemma-2-9b & $0.379$ & $0.364$ & $0.780$ & $+0.401$ \ci{+0.326}{+0.476} & $-0.016$ \ci{-0.069}{+0.036} \\
 & phi-4 & $0.379$ & $0.341$ & $0.866$ & $+0.486$ \ci{+0.415}{+0.563} & $-0.038$ \ci{-0.085}{+0.005} \\
 & Sonnet~5 & $0.379$ & $0.348$ & $0.702$ & $+0.322$ \ci{+0.249}{+0.395} & $-0.031$ \ci{-0.091}{+0.025} \\
\bottomrule
\end{tabular}
\end{table}


\section{Which correct answer best-of-$n$ returns}
\label{app:bonref}

\Cref{app:rlsim} asks what a reference reward would pull a policy toward. This asks the same
question of the thing that is actually deployed: best-of-$n$ selection, where a judge scores $n$
sampled candidates and the highest scorer is returned. No training and no simulation dynamics --
argmax over stored scores, on the tasks that carry both a matching and a non-matching correct
answer, so that there is something to choose between.\ph{}

\paragraph{Result.} \Cref{tab:bonref} gives the share of \emph{correct} picks that are the
particular correct answer the judge was shown. At $n=1$ all three selectors agree by construction,
which is the sanity check. As the budget grows, the oracle selector stays flat --- 0.18 on
arithmetic and 0.43 on \textsc{MACE}, the rate at which the matching answer occurs among correct
candidates --- while the referenced selector climbs monotonically to between 0.78 and 0.98. The
blind control not only fails to climb, it drifts slightly down ($-0.026$ to $-0.108$ at $n=64$).
Sampling more candidates does not dilute the effect; it sharpens it, because a larger pool offers
more chances for the shown answer to appear and be selected.

\paragraph{This explains a null we reported.} Best-of-64 accuracy did not fall on \textsc{MACE}'s
multi-answer questions, and our registered loss there was not confirmed (\cref{app:estimator}). The
reason is visible here: accuracy does not fall because the answer returned is still correct. What
changes is \emph{which} correct answer is returned. A pipeline measuring only accuracy would see
nothing wrong while its output distribution collapsed onto whichever answer its grader happened to
hold --- at $n=64$, 94\% and 95\% of the time for the two open judges, against 44\% under an exact
checker. For a benchmark that reports accuracy this is invisible; for anything that consumes the
selected outputs, such as distillation or preference data, it is not.

\begin{table}[htbp]
\centering
\footnotesize
\caption{\textbf{The more candidates you sample, the more reliably a reference-guided judge returns
the answer it was shown.} Share of correct best-of-$n$ picks that are the reference-matching correct
answer (95\% intervals, item bootstrap). \emph{Blind} is the same judge with no reference. At $n=1$
the three selectors are identical by construction. All post hoc.\ph}
\label{tab:bonref}
\setlength{\tabcolsep}{4pt}
\begin{tabular}{@{}llccccc@{}}
\toprule
Pool & Judge & Reward & $n{=}4$ & $n{=}16$ & $n{=}64$ & $n{=}64$ $-$ oracle \\
\midrule
Arithmetic & Qwen2.5-7B & oracle & $0.165$ & $0.183$ & $0.183$ & --- \\
 & & blind & $0.176$ & $0.163$ & $0.123$ & $-0.060$ \\
 & & referenced & $0.393$ & $0.578$ & $\mathbf{0.926}$ & $+0.744$ \\
 & phi-4 & oracle & $0.165$ & $0.183$ & $0.183$ & --- \\
 & & blind & $0.176$ & $0.188$ & $0.146$ & $-0.037$ \\
 & & referenced & $0.393$ & $0.788$ & $\mathbf{0.978}$ & $+0.795$ \\
\midrule
\textsc{MACE} & Qwen2.5-7B & oracle & $0.423$ & $0.432$ & $0.435$ & --- \\
 & & blind & $0.382$ & $0.302$ & $0.378$ & $-0.057$ \\
 & & referenced & $0.630$ & $0.811$ & $\mathbf{0.939}$ & $+0.505$ \\
 & phi-4 & oracle & $0.423$ & $0.432$ & $0.435$ & --- \\
 & & blind & $0.337$ & $0.340$ & $0.326$ & $-0.108$ \\
 & & referenced & $0.605$ & $0.865$ & $\mathbf{0.951}$ & $+0.517$ \\
 & Sonnet~5 & oracle & $0.423$ & $0.432$ & $0.435$ & --- \\
 & & blind & $0.402$ & $0.414$ & $0.409$ & $-0.026$ \\
 & & referenced & $0.526$ & $0.676$ & $\mathbf{0.783}$ & $+0.348$ \\
\bottomrule
\end{tabular}
\end{table}

\section{The KL-regularized optimum, in closed form}
\label{app:kltilt}

The simulation of \cref{app:rlsim} and the selection result of \cref{app:bonref} both approximate a
question that has an exact answer. Reinforcement learning from a reward model maximises
$\mathbb{E}[r] - \beta\,\mathrm{KL}(\pi \,\|\, \pi_{\text{base}})$, whose optimum is analytic:
$\pi^{*}(y) \propto \pi_{\text{base}}(y)\exp(r(y)/\beta)$. Both terms are already measured --
$\pi_{\text{base}}$ from how often each answer appears among the 64 stored samples, $r$ from the
stored judge scores -- so the whole optimisation path can be traced without training anything.
Sweeping $\beta$ from large (no optimisation) to small (argmax) walks up the KL budget.\ph{}

\paragraph{The method checks out against best-of-$n$.} Best-of-$n$ has an analytic KL of
$\log n - (n-1)/n$ nats, putting $n=16$ at 1.84. At that budget the closed form gives 0.815, 0.867
and 0.723 for Qwen2.5-7B, phi-4 and Sonnet~5 on \textsc{MACE}; best-of-16 measured 0.811, 0.865 and
0.676 (\cref{tab:bonref}). An analytic construction and an empirical selector agree to within two
points, which is the evidence that neither is an artifact of its own machinery.

\paragraph{What the path shows.} \Cref{tab:kltilt} gives it. We expected the reward-overoptimisation
shape of \citet{gao2022scaling}: proxy reward climbing while true reward peaks and falls, with the
reference collapse hidden in the gap. That is not what happens. Over the whole range we can reach,
accuracy climbs \emph{monotonically} alongside the collapse. At a KL budget of about 1 nat --
unremarkable for a deployed run -- Qwen2.5-7B on arithmetic goes from 46.1\% to 71.6\% accuracy
while the share of its correct mass sitting on the single displayed answer goes from 19.6\% to
77.0\%. On \textsc{MACE}, Sonnet~5 reaches 99.1\% accuracy at 1.6 nats with 69.3\% concentration.

This is worse than the overoptimisation story, not better. Overoptimisation at least eventually
announces itself: true reward turns over, and a team watching it sees something. Here the metric a
team actually monitors improves by 24 to 38 points while the answer distribution narrows onto one
arbitrary member of the correct set. There is no turnover to detect, because the reference-guided
reward is genuinely informative about correctness \emph{and} about agreement at the same time, and
only the first is being measured.

The limit of \cref{app:rlsim} still applies: $\pi_{\text{base}}$ ranges over answers that were
sampled, so this describes how optimisation pressure redistributes existing mass, not what an
exploring policy would discover. We report it as the sharpest available statement of the mechanism
and continue to claim no demonstrated harm to a trained policy.

\begin{table}[htbp]
\centering
\footnotesize
\caption{\textbf{Accuracy and reference concentration rise together; nothing warns you.} The
KL-regularized optimum $\pi^{*} \propto \pi_{\text{base}}\exp(r/\beta)$ computed in closed form from
stored scores. \emph{Proxy}: mean judge score under $\pi^{*}$. \emph{True}: oracle accuracy.
\emph{P(ref\,$\vert$\,correct)}: share of correct mass on the displayed answer. All post hoc.\ph}
\label{tab:kltilt}
\setlength{\tabcolsep}{4pt}
\begin{tabular}{@{}llccccc@{}}
\toprule
Pool & Judge & KL (nats) & Proxy & True accuracy & P(ref\,$\vert$\,correct) \\
\midrule
Arithmetic & Qwen2.5-7B & $0.00$ (base) & $4.23$ & $0.461$ & $0.196$ \\
 & & $0.72$ & $7.07$ & $0.622$ & $0.551$ \\
 & & $1.97$ & $8.83$ & $0.716$ & $0.770$ \\
 & & $3.30$ & $9.60$ & $0.736$ & $0.929$ \\
\cmidrule(l){2-6}
 & phi-4 & $0.00$ (base) & $5.83$ & $0.461$ & $0.196$ \\
 & & $0.54$ & $8.18$ & $0.605$ & $0.464$ \\
 & & $1.34$ & $9.31$ & $0.690$ & $0.654$ \\
 & & $2.72$ & $9.99$ & $0.899$ & $0.896$ \\
\midrule
\textsc{MACE} & Qwen2.5-7B & $0.00$ (base) & $3.47$ & $0.614$ & $0.381$ \\
 & & $0.91$ & $7.37$ & $0.795$ & $0.704$ \\
 & & $1.76$ & $8.62$ & $0.850$ & $0.815$ \\
 & & $2.93$ & $9.11$ & $0.870$ & $0.910$ \\
\cmidrule(l){2-6}
 & phi-4 & $0.00$ (base) & $4.79$ & $0.614$ & $0.381$ \\
 & & $0.75$ & $7.90$ & $0.799$ & $0.709$ \\
 & & $1.65$ & $9.21$ & $0.878$ & $0.867$ \\
 & & $2.55$ & $9.71$ & $0.928$ & $0.940$ \\
\cmidrule(l){2-6}
 & Sonnet~5 & $0.00$ (base) & $3.94$ & $0.614$ & $0.381$ \\
 & & $1.08$ & $8.86$ & $0.953$ & $0.628$ \\
 & & $1.62$ & $9.68$ & $0.991$ & $0.693$ \\
 & & $1.85$ & $9.83$ & $1.000$ & $0.723$ \\
\bottomrule
\end{tabular}
\end{table}